% \iffalse meta-comment
%
% Copyright (C) 2017 Universite Laval
%
% This file may be distributed and/or modified under the conditions
% of the LaTeX Project Public License, either version 1.3c of this
% license or (at your option) any later version. The latest version
% of this license is in:
%
%   http://www.latex-project.org/lppl.txt
%
% and version 1.3c or later is part of all distributions of LaTeX
% version 2006/05/20 or later.
%
% This work has the LPPL maintenance status `maintained'.
%
% The Current Maintainer of this work is Universite Laval
% <ulthese-dev@bibl.ulaval.ca>.
%
% This work consists of the files ulthese.dtx and ulthese.ins
% and the derived files listed in the README.md file.
%
% \fi
%
% \iffalse
%<*dtx>
\ProvidesFile{ulthese.dtx}
%</dtx>
%<class>\NeedsTeXFormat{LaTeX2e}[2009/09/24]
%<class>\ProvidesClass{ulthese}%
%<*class>
  [2017/06/01 v4.4 Universite Laval thesis and memoir class]
%</class>
%<*driver>
\documentclass[11pt,english,french]{ltxdoc}
  \usepackage[utf8]{inputenc}
  \usepackage[T1]{fontenc}
  \usepackage{natbib}
  \usepackage{babel}
  \usepackage[autolanguage]{numprint}
  \usepackage{microtype}
  \usepackage[scaled=0.92]{helvet}
  \usepackage[scaled=1.02]{inconsolata}
  \usepackage[sc]{mathpazo}
  \usepackage{fontawesome}
  \usepackage{metalogo}
  \usepackage{tabularx,booktabs}
  \DisableCrossrefs
  \CodelineNumbered
  \RecordChanges
  \GlossaryPrologue{\section*{Historique des versions}%
    \addcontentsline{toc}{section}{Historique des versions}}

  %% Couleurs
  \usepackage{xcolor}
  \definecolor{link}{rgb}{0,0.4,0.6}   % ~RoyalBlue de dvips
  \definecolor{url}{rgb}{0.6,0,0}      % rouge foncé
  \definecolor{citation}{rgb}{0,0.5,0} % vert foncé

  %% Liste description alignée à gauche
  \usepackage{enumitem}
  \setlist[description]{leftmargin=*,align=left}

  %% Environnement pour remarques
  \usepackage{amsthm}
  \theoremstyle{definition}
  \newtheorem*{rem}{Remarque}

  %% Hyperliens
  \usepackage{hyperref}
  \hypersetup{%
    pdfauthor = {Faculté des études supérieures et postdoctorales},
    pdftitle = {Guide de l'utilisateur de la classe ulthese},
    colorlinks = {true},
    linktocpage = {true},
    urlcolor = {url},
    linkcolor = {link},
    citecolor = {citation},
    pdfpagemode = {UseOutlines},
    pdfstartview = {Fit},
    bookmarksopen = {true},
    bookmarksnumbered = {true},
    bookmarksdepth = {subsubsection}}
  \addto\extrasfrench{%
    \def\appendixautorefname{annexe}%
    \def\tableautorefname{tableau}%
    \def\subsectionautorefname{section}%
  }

  %% Paramétrage de babel
  \frenchbsetup{%
    StandardItemizeEnv=true,       % listes standards
    ThinSpaceInFrenchNumbers=true, % espace fine dans les nombres
    og=«, fg=»                     % « et » sont les guillemets
  }
  \def\frenchtablename{{\scshape Tab.}}

  %% Commandes de mise en forme spéciales
  \newcommand{\class}[1]{\textsf{#1}}
  \newcommand{\pkg}[1]{\textbf{#1}}

  %% Commandes pour les liens vers la documentation
  \newcommand{\link}[2]{\href{#1}{#2~\raisebox{-0.2ex}{\faExternalLink}}}
  \newcommand{\doc}[3][documentation]{%
    \link{#3}{#1}\marginpar{\hfill\faBookmark~\texttt{#2}}}

\begin{document}
  \DocInput{ulthese.dtx}
\end{document}
%</driver>
% \fi
% \CheckSum{0}
% \DoNotIndex{\',\^,\`,\ ,\ae}
% \DoNotIndex{\RequirePackage,\ExecuteOptions,\ProcessOptions}
% \DoNotIndex{\newcommand,\newcommand*}
% \DoNotIndex{\setlength}
% \changes{4.4}{2017-05-16}{Prise en charge des bibliographies
%   multiples pour les thèses et mémoires par articles.}
% \changes{4.4}{2017-05-16}{Nouveau gabarit pour une thèse par
%   articles.}
% \changes{4.4}{2017-05-16}{Documentation révisée, principalement ce
%   qui touche à la préparation de la bibliographie.}
% \changes{4.3}{2017-03-01}{Nouvelles options pour les titres de
%   docteur en musique, maître en architecture, maître en
%   ergothérapie, maître en physiothérapie et maître en
%   psychoéducation.}
% \changes{4.3}{2017-03-01}{Modifications à la composition des sigles de
%   grades.}
% \changes{4.3}{2017-03-01}{Vérification de la compatibilité entre le
%   grade et les options multifacultaire, cotutelle, bidiplomation et
%   extension.}
% \changes{4.3}{2017-03-01}{Améliorations à la documentation.}
% \changes{4.2}{2016-03-29}{Ajout de l'option \texttt{examen}. Nouvelles
%   options de \pkg{babel} dans les gabarits.}
% \changes{4.1}{2016-02-13}{Paquetage \pkg{geometry} déclaré incompatible
%   avec la classe.}
% \changes{4.0}{2015-06-12}{Nouvelles règles de présentation matérielle
%   de la FESP: recto seulement, page frontispice.}
% \changes{3.1}{2014-05-23}{Prise en charge de la maîtrise en bidiplomation}
% \changes{3.0a}{2014-03-24}{Modifications et corrections à la
%   documentation, notamment relativement à la configuration de \pkg{natbib}}
% \changes{3.0}{2014-01-06}{Déclaration du grade en option de la
%   classe.}
% \changes{3.0}{2014-01-06}{Moteur {\XeLaTeX} supporté}
% \changes{3.0}{2014-01-06}{Ajout de l'option \texttt{nobabel}}
% \changes{2.1}{2013-01-16}{Utilisation transparente de la police
%   Helvetica pour la page de titre}
% \changes{2.0}{2013-01-13}{Traitement automatique des longs titres}
% \changes{1.0b}{2012-11-11}{Ajouts et corrections mineures dans la
%   documentation}
% \changes{1.0a}{2012-10-17}{Précisions dans la documentation}
% \changes{1.0}{2012-09-30}{Version initiale}
% \GetFileInfo{ulthese.dtx}
% \title{\class{ulthese}: la classe pour les thèses et mémoires de
%   l'Université Laval\thanks{Ce document décrit la classe
%   \class{ulthese}~\fileversion, datée du \filedate.}}
% \author{Faculté des études supérieures et postdoctorales\thanks{%
%   Cette classe et sa documentation ont été rédigées par Vincent
%   Goulet~(Faculté des sciences et de génie) avec la collaboration de
%   Koassi D'Almeida~(Faculté des études supérieures et postdoctorales) et
%   de Pierre Lasou~(Bibliothèque).}}
% \date{}
% \maketitle
%
% \section{Introduction}
%
% La classe \class{ulthese} permet de composer avec {\LaTeX} ou
% {\XeLaTeX} des thèses et mémoires immédiatement conformes aux règles
% générales de présentation matérielle de la Faculté des études
% supérieures et postdoctorales (FESP) de l'Université Laval. Ces
% règles définissent principalement la présentation des pages de titre
% des thèses et mémoires ainsi que la disposition du texte sur la
% page. La classe en elle-même est donc relativement simple.
%
% Cependant, \class{ulthese} est basée sur la classe \class{memoir}
% \citep{memoir}, une extension de la classe standard \class{book}
% facilitant à plusieurs égards la préparation de documents d'allure
% professionnelle dans {\LaTeX}. La classe \class{memoir} incorpore
% d'office plus de 30 des paquetages les plus populaires\footnote{%
% Consulter la section~18.24 de la documentation de \class{memoir}
% pour la liste ou encore le journal de la compilation (\emph{log})
% d'un document utilisant la classe \class{ulthese}.}. %
% L'intégralité des fonctionnalités de \class{memoir} se retrouve donc
% dans \class{ulthese}.
%
% La classe \class{memoir} fait partie des distributions {\LaTeX}
% modernes; elle devrait donc être installée et disponible sur votre
% système. Elle est livrée avec une %
% \doc{memman}{http://texdoc.net/pkg/memoir}\footnote{%
% Première occurence d'une convention de ce document quand il s'agit
% de documentation d'un paquetage: un hyperlien mène vers une version
% en ligne dans le site \link{http://texdoc.net}{TeXdoc Online} et on
% trouve dans la marge le nom du fichier correspondant (sans
% l'extension \texttt{.pdf}) sur un système doté d'une installation de
% {\TeX}~Live.} %
% exhaustive: le guide de l'utilisateur fait près de 600~pages!
% N'hésitez pas à vous y référer pour réaliser une mise en page
% particulière.
%
% L'autre compagnon naturel de la présente documentation est
% \doc[Rédaction avec {\LaTeX}]{formation-latex-ul}{http://texdoc.net/pkg/formation-latex-ul} ^^A
% \citep{formation-latex-ul}, la formation {\LaTeX} de l'Université Laval.
% Vous y trouverez des informations additionnelles sur l'utilisation des
% classes \class{ulthese} et \class{memoir}.
%
% \section{Démarrage rapide (pour les impatients)}
% \label{sec:utilisation:rapide}
%
% La classe est livrée avec les distributions {\TeX}~Live, Mac{\TeX} et
% MiK{\TeX}. Si vous utilisez l'une ou l'autre de ces distributions et
% qu'elle est à jour\footnote{%
%   Validez à l'aide de l'assistant de mise à jour de votre
%   distribution.}, %
% vous devriez pouvoir utiliser \class{ulthese} sans autre
% intervention.
%
% Il est recommandé de segmenter tout document d'une certaine ampleur
% dans des fichiers |.tex| distincts pour chaque partie ---
% habituellement un fichier par chapitre. Le document complet est
% composé à l'aide d'un fichier maître qui contient le préambule
% {\LaTeX} et un ensemble de commandes \cmd{\include} pour réunir les
% parties dans un tout.
%
% La classe \class{ulthese} est livrée avec un ensemble de gabarits sur
% lesquels se baser pour:
% \begin{itemize}
% \item les fichiers maîtres de divers types de thèses et de mémoires
%   (standard, par articles, sur mesure, en cotutelle, en
%   bidiplomation, en extension, etc.);
% \item les fichiers des parties les plus usuelles (résumés français
%   et anglais, avant-propos, introduction, chapitres, conclusion,
%   etc.).
% \end{itemize}
% Les noms des fichiers devraient permettre de facilement identifier
% leur contenu (c'est là une bonne pratique: le nom |rappels.tex|
% parle de lui-même et résiste mieux aux changements à l'ordre des
% chapitres que |chapitre1.tex|).
%
% Dans {\TeX}~Live, les gabarits sont classés avec la documentation de
% la classe. Pour les utiliser, copiez les fichiers appropriés dans votre
% dossier de travail.
%
% Pour débuter la rédaction, renommez le gabarit de document maître
% approprié d'après votre numéro de dossier. Par exemple, l'étudiante
% dont le numéro de dossier est 900352789 et qui entame la rédaction
% d'une thèse multifacultaire renommera le fichier
% \begin{quote}
%   |gabarit-doctorat-multifacultaire.tex|
% \end{quote}
% en
% \begin{quote}
%   |900352789.tex|.
% \end{quote}
%
% Les gabarits comportent des commentaires succincts pour vous guider
% dans la préparation de votre document. La \autoref{sec:gabarits}
% fournit de plus amples informations.
%
% \section{Installation}
%
% Cette section explique comment installer la classe \class{ulthese}
% si elle n'est pas disponible sur votre système ou si la version
% n'est pas à jour.
%
% La classe est distribuée sous forme d'une archive |ulthese.zip| via
% le réseau de sites \emph{Comprehensive {\TeX} Archive Network} (CTAN):
% \begin{quote}
%   \url{https://www.ctan.org/pkg/ulthese}
% \end{quote}
%
% L'installation de la classe consiste à créer le fichier
% |ulthese.cls| et plusieurs gabarits |.tex| à partir du code source
% documenté se trouvant dans le fichier |ulthese.dtx|. Il est
% recommandé de simplement créer ces fichiers dans le dossier de
% travail de la thèse ou du mémoire.
%
% Pour procéder à l'installation, décompressez l'archive |ulthese.zip|
% dans votre dossier de travail, puis compilez avec {\LaTeX} le fichier
% |ulthese.ins| en exécutant
% \begin{quote}
%   |latex ulthese.ins|
% \end{quote}
% depuis une invite de commande. Il est aussi possible d'ouvrir le
% fichier |ulthese.ins| dans votre éditeur de texte favori et de
% lancer depuis celui-ci la compilation avec {\LaTeX}, pdf{\LaTeX},
% {\XeLaTeX} ou un autre moteur {\TeX}.
%
% \section{Utilisation}
%
% La classe est compatible avec les moteurs {\LaTeX} traditionnels
% ainsi qu'avec le plus récent moteur {\XeLaTeX}.
%
% On charge la classe avec la commande
% \begin{quote}
% |\documentclass|\oarg{options}|{ulthese}|
% \end{quote}
% Les marges, l'interligne et la numérotation des pages sont adaptées
% aux règles de présentation matérielle de la FESP. Les sections
% suivantes décrivent les options et les commandes définies par la
% classe.
%
% \subsection{Options de la classe}
% \label{sec:utilisation:options}
%
% Cette section passe en revue les \meta{options} que l'on peut
% spécifier au chargement de la classe. Les commandes mentionnées
% ci-dessous font quant à elles l'objet de la \autoref{sec:commandes}.
%
% \begin{DescribeMacro}{PhD,MSc,MA, ...}
%   Identifie le type de grade; consulter le \autoref{tab:grades}
%   pour la liste complète des options et les grades correspondants.
%   La déclaration d'un type de grade est obligatoire.
% \end{DescribeMacro}
%
% \begin{table}[t]
%   \centering
%   \begin{tabularx}{0.85\linewidth}{lX}
%     \toprule
%     Option & Nom du grade (sigle) \\
%     \midrule
%     |LLD|
%       & Docteur en droit (LL.~D.) \\
%     |DMus|
%       & Docteur en musique (D.~Mus.) \\
%     |DPsy|
%       & Docteur en psychologie (D.~Psy.) \\
%     |DThP|
%       & Docteur en théologie pratique (D.~Th.~P.) \\
%     |PhD|
%       & Philosophi{\ae} doctor (Ph.~D.) \\
%     \addlinespace[6pt]
%     |MATDR|
%       & Maître en aménagement du territoire et développement régional
%       (M.ATDR) \\
%     |MArch|
%       & Maître en architecture
%       (M.~Arch.) \\
%     |MA|
%       & Maître ès arts (M.A.) \\
%     |LLM|
%       & Maître en droit (LL.~M.) \\
%     |MErg|
%       & Maître en ergothérapie (M.~Erg.) \\
%     |MMus|
%       & Maître en musique (M.~Mus.) \\
%     |MPht|
%       & Maître en physiothérapie (M.~Pht.) \\
%     |MSc|
%       & Maître ès sciences (M.~Sc.) \\
%     |MScGeogr|
%       & Maître en sciences géographiques (M.~Sc.~géogr.) \\
%     |MServSoc|
%       & Maître en service social (M.~Serv.~soc.) \\
%     |MPsEd|
%       & Maître en psychoéducation (M.~Ps.~éd.) \\
%     \bottomrule
%   \end{tabularx}
%   \caption{Options de la classe pour la déclaration du grade et libellés
%     correspondants}
%   \label{tab:grades}
% \end{table}
%
% \begin{DescribeMacro}{multifacultaire}
%   Identifie une thèse multifacultaire. Valide uniquement avec un
%   grade de doctorat. Cette déclaration requiert ensuite d'utiliser
%   la commande \cmd{\faculteUL}.
% \end{DescribeMacro}
%
% \begin{DescribeMacro}{cotutelle}
%   Identifie une thèse effectuée en cotutelle avec une autre
%   université. Valide uniquement avec un grade de doctorat. Cette
%   déclaration requiert ensuite d'utiliser les commandes
%   \cmd{\univcotutelle} et \cmd{\gradecotutelle}.
% \end{DescribeMacro}
%
% \begin{DescribeMacro}{bidiplomation}
%   Identifie un mémoire effectué en bidiplomation avec une autre
%   université. Valide uniquement avec un grade de maîtrise. Cette
%   déclaration requiert ensuite d'utiliser les commandes
%   \cmd{\univbidiplomation} et \cmd{\gradebidiplomation}.
% \end{DescribeMacro}
%
% \begin{DescribeMacro}{extensionUdeS}
%   Identifie une thèse réalisée en extension à l'Université de
%   Sherbrooke. Valide uniquement avec un grade de doctorat. Cette
%   déclaration requiert ensuite d'utiliser les commandes
%   \cmd{\faculteUL} et \cmd{\faculteUdeS}.
% \end{DescribeMacro}
%
% \begin{DescribeMacro}{extensionUQO}
%   Identifie une thèse réalisée en extension à l'Université du
%   Québec en Outaouais (UQO). Valide uniquement avec un grade de doctorat.
%   Cette déclaration requiert ensuite d'utiliser les commandes
%   \cmd{\faculteUL} et \cmd{\faculteUQO}.
% \end{DescribeMacro}
%
% \begin{DescribeMacro}{extensionUQAC}
%   Identifie un mémoire réalisé en extension à l'Université du Québec
%   à Chicoutimi (UQAC). Valide uniquement avec un grade de maîtrise.
%   Cette déclaration requiert ensuite d'utiliser les commandes
%   \cmd{\faculteUL} et \cmd{\faculteUQAC}.
% \end{DescribeMacro}
%
% \begin{DescribeMacro}{examen}
%   Identifie un examen de doctorat. Cette option permet d'utiliser la
%   classe pour la rédaction d'un examen de doctorat respectant les
%   règles de présentation matérielle de la FESP. Elle a pour effet de
%   changer l'appellation «Thèse» sur la page couverture pour «Examen
%   de doctorat». Elle supprime également la page frontispice.
%   L'option n'est compatible qu'avec l'une des options de grade de
%   doctorat.
% \end{DescribeMacro}
%
% \begin{DescribeMacro}{10pt,11pt,12pt}
%   Sélectionne une taille de police de 10, 11 ou 12~points. Par
%   défaut la classe utilise une police de 11~points. Ces options
%   n'ont aucun effet sur la taille des polices des pages de titre.
% \end{DescribeMacro}
%
% \changes{4.4}{2017-05-15}{Ajout de l'option \texttt{bibchapitre}}
% \changes{4.4}{2017-05-15}{Ajout de l'option \texttt{bibsection}}
% \begin{DescribeMacro}{bibchapitre}
%   \begin{DescribeMacro}{bibsection}
%     Permettent de composer une bibliographie distincte après chaque
%     fichier inséré dans le document avec \cmd{\include},
%     habituellement un chapitre; voir la
%     \autoref{sec:bibliographie:multiples}. Ces options sont
%     principalement utiles pour les thèses et mémoires par articles.
%     La bibliographie se présente sous la forme d'un chapitre non
%     numéroté avec |bibchapitre|, et sous la forme d'une section
%     numérotée avec |bibsection|.
%   \end{DescribeMacro}
% \end{DescribeMacro}
%
% \begin{DescribeMacro}{nonatbib}
%   Empêche le chargement du paquetage \pkg{natbib}. La paquetage est
%   normalement chargé par la classe; voir la
%   \autoref{sec:bibliographie}. L'option |nonatbib| permet d'empêcher
%   le chargement pour modifier les options du paquetage ou en cas de
%   conflit avec un autre paquetage de mise en forme de la
%   bibliographie.
% \end{DescribeMacro}
%
% \begin{DescribeMacro}{nobabel}
%   Empêche le chargement du paquetage \pkg{babel}. La classe utilise
%   par défaut ce paquetage pour le traitement des langues dans le
%   document; voir la \autoref{sec:langues}. L'option |nobabel| permet
%   d'empêcher son chargement si un autre paquetage devait être
%   utilisé --- on pense ici principalement à \pkg{polyglossia} pour
%   un document produit avec le moteur {\XeLaTeX}.
% \end{DescribeMacro}
%
% \begin{DescribeMacro}{english, french, ...}
%   Déclare les langues utilisées dans le document. Ces options sont
%   transférées au paquetage \pkg{babel} (dans la mesure où |nobabel|
%   n'est pas spécifié, bien entendu). Le libellé des langues devrait
%   donc correspondre aux options de \pkg{babel}. La dernière langue
%   spécifiée est la langue active par défaut dans le document.
% \end{DescribeMacro}
%
% Toute autre option sera passée à la classe \class{memoir} dont,
% entre autres, le format du papier. Le format lettre nord-américain
% (option |letterpaper|) est utilisé par défaut. Si votre thèse doit
% être imprimée en format international A4, utilisez l'option
% |a4paper|. La classe \class{memoir} est toujours chargée avec
% l'option |oneside|.
%
% \subsection{Commandes de la classe}
% \label{sec:commandes}
%
% La classe \class{ulthese} définit quelques nouvelles commandes
% servant principalement à créer les pages de titre et des éléments
% des pages liminaires. On trouvera un sommaire des commandes au
% \autoref{tab:commandes} et leurs descriptions détaillées ci-dessous.
%
% \begin{table}[t]
%   \centering
%   \begin{tabularx}{\linewidth}{lX}
%     \toprule
%     Commande & Usage \\
%     \midrule
%     \cmd{\titre}$^\star$ & titre principal du document \\
%     \cmd{\soustitre} & sous-titre du document \\
%     \cmd{\auteur}$^\star$ & nom complet de l'auteur \\
%     \cmd{\annee}$^\star$ & année du dépôt final \\
%     \cmd{\programme}$^\star$ & nom officiel du programme d'études \\
%     \cmd{\direction}$^\star$ & nom du directeur (directrice) de recherche \\
%     \cmd{\codirection} & noms des codirecteurs de recherche \\
%     \cmd{\univcotutelle} & université de cotutelle \\
%     \cmd{\gradecotutelle} & grade conféré par l'université de cotutelle \\
%     \cmd{\univbidiplomation} & université de bidiplomation \\
%     \cmd{\gradebidiplomation} & grade conféré par l'université de bidiplomation \\
%     \cmd{\faculteUL} & noms des facultés de l'Université Laval \\
%     \cmd{\faculteUdeS} & nom de la faculté de l'Université de Sherbrooke \\
%     \cmd{\faculteUQO} & nom de la faculté de l'UQO \\
%     \cmd{\faculteUQAC} & nom de la faculté de l'UQAC \\
%     \cmd{\pagestitre}$^\star$ & production des pages de titre et frontispice \\
%     \cmd{\dedicace} & dédicace du document \\
%     \cmd{\epigraphe} & épigraphe du document \\
%     \bottomrule
%   \end{tabularx}
%   \caption{Sommaire des commandes de la classe \class{ulthese}.
%     Celles marquées d'une étoile $^\star$ sont obligatoires.}
%   \label{tab:commandes}
% \end{table}
%
% \begin{DescribeMacro}{\titre}
%   Titre principal de la thèse ou du mémoire. Ne pas utiliser la
%   commande \cmd{\title} de {\LaTeX} pour ce faire.
%
%   Un titre très long devra être coupé manuellement avec |\\| ou
%   \cmd{\newline}. Par exemple, la déclaration d'un titre d'une seule
%   ligne est:
%   \begin{quote}
%     |\titre{Ceci est un titre d'une seule ligne}|
%   \end{quote}
%   Pour un titre de deux lignes, utiliser:
%   \begin{quote}
%     |\titre{Ceci est la première ligne d'un long titre \\| \\
%     |       et ceci est la seconde}|
%   \end{quote}
% \end{DescribeMacro}
%
% \begin{DescribeMacro}{\soustitre}
%   Sous-titre de la thèse ou du mémoire, le cas échéant. Les remarques
%   sur un long titre principal s'appliquent également au sous-titre.
% \end{DescribeMacro}
%
% \begin{DescribeMacro}{\auteur}
%   Nom complet de l'auteur de la thèse ou du mémoire, sous la forme
%   |Prénom Nom| avec seulement des majuscules initiales. Ne pas
%   utiliser la commande \cmd{\author} de {\LaTeX} pour le nom de
%   l'auteur.
% \end{DescribeMacro}
%
% \begin{DescribeMacro}{\annee}
%   Année du dépôt final de la thèse ou du mémoire.
% \end{DescribeMacro}
%
% \begin{DescribeMacro}{\programme}
%   Nom complet officiel du programme d'études comme «Doctorat en
%   informatique» ou «Maîtrise en mathématiques». Si le programme
%   comporte une majeure, séparer sa mention de celle du programme
%   principal par un tiret demi-quadratin (obtenu avec |--|).
% \end{DescribeMacro}
%
% \begin{DescribeMacro}{\direction}
%   Nom complet du directeur ou de la directrice de recherche, sous la
%   forme |Prénom Nom| avec seulement des majuscules initiales, suivi
%   d'une virgule et de la mention «directeur de recherche» ou
%   «directrice de recherche».
%
%   Les thèses en cotutelle comportent un directeur ou une directrice
%   de recherche et un directeur ou une directrice de cotutelle.
%   Séparer chaque mention par |\\|, comme ceci:
%   \begin{quote}
%     |\direction{Prénom Nom, directrice de recherche \\| \\
%     |           Prénom Nom, directeur de cotutelle}|
%   \end{quote}
%
%   De même, les maîtrises en bidiplomation comptent deux directeurs
%   ou directrices de recherche. Séparer chaque mention comme
%   mentionné ci-dessus.
% \end{DescribeMacro}
%
% \begin{DescribeMacro}{\codirection}
%   Comme la commande \cmd{\direction}, mais pour le ou les
%   codirecteurs de recherche, s'il y a lieu. Lorsqu'il y a plus d'un
%   codirecteur de recherche, séparer chaque mention par |\\|, comme
%   ceci:
%   \begin{quote}
%     |\codirection{Prénom Nom, directrice de recherche \\| \\
%     |             Prénom Nom, directeur de recherche}|
%   \end{quote}
% \end{DescribeMacro}
%
% \begin{DescribeMacro}{\univcotutelle}
%   Nom, ville et pays de l'université de cotutelle, déclarés sous la
%   forme
%   \begin{quote}
%     |\univcotutelle{Nom de l'université \\ Ville, Pays}|
%   \end{quote}
%   Cette commande prend effet seulement lorsque la classe est chargée
%   avec l'option |cotutelle|.
% \end{DescribeMacro}
%
% \begin{DescribeMacro}{\gradecotutelle}
%   Grade conféré par l'université de cotutelle, déclaré
%   sous la forme
%   \begin{quote}
%     |\gradecotutelle{Nom du grade (sigle)}|
%   \end{quote}
%   Cette commande prend effet seulement lorsque la classe est chargée
%   avec l'option |cotutelle|.
% \end{DescribeMacro}
%
% \begin{DescribeMacro}{\univbidiplomation}
%   Nom, ville et pays de l'université de bidiplomation, déclarés sous
%   la forme
%   \begin{quote}
%     |\univbidiplomation{Nom de l'université \\ Ville, Pays}|
%   \end{quote}
%   Cette commande prend effet seulement lorsque la classe est chargée
%   avec l'option |bidiplomation|.
% \end{DescribeMacro}
%
% \begin{DescribeMacro}{\gradebidiplomation}
%   Grade conféré par l'université de bidiplomation, déclaré sous la
%   forme
%   \begin{quote}
%     |\gradebidiplomation{Nom du grade (sigle)}|
%   \end{quote}
%   Cette commande prend effet seulement lorsque la classe est chargée
%   avec l'option |bidiplomation|.
% \end{DescribeMacro}
%
% \begin{DescribeMacro}{\faculteUL}
%   La commande a deux usages:
%   \begin{enumerate}
%   \item noms des facultés pour les thèses et mémoires
%     multifacultaires, séparés par des commandes |\\|;
%   \item nom de la faculté de l'Université Laval où sont réalisés les
%     thèses et mémoires en extension à l'Université de Sherbrooke, à
%     l'UQO ou à l'UQAC.
%   \end{enumerate}
%   Cette commande prend effet seulement lorsque la classe est chargée
%   avec l'une ou l'autre des options |multifacultaire|,
%   |extensionUdeS|, |extensionUQO| ou |extensionUQAC|.
% \end{DescribeMacro}
%
% \begin{DescribeMacro}{\faculteUdeS}
%   Nom de la faculté de l'Université de Sherbrooke hébergeant la
%   thèse en extension. Cette commande prend effet seulement lorsque
%   la classe est chargée avec l'option |extensionUdeS|.
% \end{DescribeMacro}
%
% \begin{DescribeMacro}{\faculteUQO}
%   Nom de la faculté de l'Université du Québec en Outaouais
%   hébergeant la thèse en extension. Cette commande prend effet
%   seulement lorsque la classe est chargée avec l'option
%   |extensionUQO|.
% \end{DescribeMacro}
%
% \begin{DescribeMacro}{\faculteUQAC}
%   Nom de la faculté de l'Université du Québec à Chicoutimi
%   hébergeant le mémoire en extension. Cette commande prend effet
%   seulement lorsque la classe est chargée avec l'option
%   |extensionUQAC|.
% \end{DescribeMacro}
%
% \begin{DescribeMacro}{\pagestitre}
%   Création de la page de titre et de la page frontispice. Ne pas
%   utiliser la commande \cmd{\pagetitle} de {\LaTeX} pour ce faire.
%   De toutes les commandes ci-dessus, c'est la seule qui doit se
%   trouver dans le corps du document plutôt que dans le préambule.
% \end{DescribeMacro}
%
% \begin{DescribeMacro}{\dedicace}
%   Ajout d'une dédicace («À mes parents», «À Camille») à la thèse ou
%   au mémoire. La dédicace est disposée seule sur une page liminaire,
%   à une dizaine de lignes de la marge du haut et alignée à droite.
%   Par défaut, elle est composée en italique.
% \end{DescribeMacro}
%
% \begin{DescribeMacro}{\epigraphe}
%   Ajout d'une épigraphe au début du document. Comme la dédicace,
%   l'épigraphe est disposée seule sur une page liminaire, à une
%   dizaine de lignes de la marge du haut et alignée à droite. La
%   commande accepte deux arguments, soit le texte de la citation et
%   son auteur ou la source, dans l'ordre.
%
%   Pour ajouter une épigraphe au début d'un ou plusieurs chapitres,
%   utiliser directement la commande \cmd{\epigraph} de
%   \class{memoir}, sur laquelle \cmd{\dedicace} et \cmd{\epigraphe}
%   sont d'ailleurs basées.
% \end{DescribeMacro}
%
% \subsection{Citations}
% \label{sec:citations}
%
% {\LaTeX} offre deux environnements pour les citations dans le
% texte: |quote| et |quotation|.
%
% \begin{DescribeEnv}{quote}
%   L'environnement |quote| sert pour les citations «courtes»,
%   quelques lignes au plus. Dans la classe, le texte est alors placé
%   en retrait des marges normales de 10~mm à gauche et à droite.
% \end{DescribeEnv}
%
% \begin{DescribeEnv}{quotation}
%   L'environnement |quotation|, quant à lui, doit être utilisé pour
%   les citations «longues», celles qui peuvent s'étendre sur plus de
%   cinq lignes ou, surtout, plus d'un paragraphe. Dans la classe, le
%   texte est alors toujours placé en retrait de 10~mm, mais également
%   à interligne simple. De plus, un espace vertical sépare les
%   paragraphes, le cas échéant, afin de bien les distinguer les uns
%   des autres.
% \end{DescribeEnv}
%
% \subsection{Interligne}
% \label{sec:interligne}
%
% \begin{DescribeMacro}{\OnehalfSpacing}
%   L'espacement d'un interligne et demi utilisé dans la classe est
%   obtenu avec la commande \cmd{\OnehalfSpacing} de \class{memoir}.
%   L'interligne simple est automatiquement rétabli pour les pages
%   de titre, la table des matières, la liste des tableaux, la liste
%   des figures et les longues citations (\autoref{sec:citations}).
% \end{DescribeMacro}
%
% \begin{DescribeMacro}{\SingleSpacing}
%   Si ce devait être nécessaire ailleurs dans le document, la
%   commande \cmd{\SingleSpacing} permet de passer à l'interligne simple.
% \end{DescribeMacro}
%
%
% \subsection{Autres paquetages chargés}
% \label{sec:paquetages}
%
% Outre \class{memoir}, la classe \class{ulthese} charge quelques
% paquetages qui peuvent aussi vous être utiles. Il n'est donc pas
% nécessaire de charger de nouveau les paquetages suivants:
% \begin{description}
% \item[\normalfont\pkg{babel}] \citep{babel} gestion des documents
%   rédigés dans une ou plusieurs langues autres que l'anglais (si
%   l'option |nobabel| de la classe est absente; voir aussi la
%   \autoref{sec:langues});
% \item[\normalfont\pkg{numprint}] \citep{numprint} requis par la
%   commande \cmd{\nombre} de \pkg{babel}; le paquetage est donc
%   chargé uniquement si \pkg{babel} l'est. Permet de composer
%   automatiquement des nombres avec un séparateur toutes les trois
%   positions (une espace en français);
% \item[\normalfont\pkg{natbib}] \citep{natbib} gestion de la
%   bibliographie (si l'option |nonatbib| de la classe est absente;
%   voir aussi la \autoref{sec:bibliographie});
% \item[\normalfont\pkg{chapterbib}] \citep{chapterbib} gestion de
%   bibliographies multiples (si l'une ou l'autre des options
%   |bibchapitre| ou |bibsection| est spécifiée;
%   voir aussi la \autoref{sec:bibliographie:multiples});
% \item[\normalfont\pkg{fontspec}] \citep{fontspec} gestion des
%   polices OpenType sous {\XeLaTeX} (chargé avec ce moteur
%   seulement);
% \item[\normalfont\pkg{unicode-math}] \citep{unicode-math} gestion
%   des polices mathématiques {\XeLaTeX} (chargé avec ce moteur
%   seulement);
% \item[\normalfont\pkg{graphicx}] \citep{graphicx} insertion et
%   manipulation de graphiques;
% \item[\normalfont\pkg{xcolor}] \citep{xcolor} gestion des couleurs
%   dans le document;
% \item[\normalfont\pkg{textcomp}] multitude de symboles spéciaux,
%   dont un beau symbole de copyright, \textcopyright.
% \end{description}
% L'\autoref{sec:meo} sur la mise en {\oe}uvre de la classe fournit
% plus de détails sur la liste des paquetages chargés et les raisons
% pour lesquelles ils sont requis dans la classe.
%
% \subsection{Paquetage incompatible}
% \label{sec:incompatible}
%
% Le paquetage \pkg{geometry} est incompatible avec la classe à cause
% de sa mauvaise interaction avec \class{memoir}. Son chargement dans
% le préambule du document cause une erreur lors de la compilation.
%
% \section{Français et autres langues}
% \label{sec:langues}
%
% Une complication additionnelle pour les auteurs rédigeant dans une
% langue autre que l'anglais consiste à adapter {\LaTeX} à leur
% langue, qu'il s'agisse des mots clés, de la typographie ou de la
% césure des mots. La solution standard à ce problème provient du
% paquetage \pkg{babel}. Celui-ci permet de combiner plusieurs
% langues dans un même document et de passer de l'une à l'autre
% facilement. Il est chargé par défaut par la classe \class{ulthese}.
%
% Aucune langue n'est spécifiée dans la classe. La plupart des auteurs
% auront recours à l'anglais et au français, ne serait-ce que pour les
% deux résumés demandés par la FESP. Les langues utilisées dans le
% document doivent être spécifiées comme options à la classe, tel que
% mentionné à la \autoref{sec:utilisation:options}. La
% \emph{dernière} langue spécifiée devient par défaut la langue active
% du document.
%
% \begin{DescribeMacro}{\selectlanguage}
%   La commande \cmd{\selectlanguage} de \pkg{babel} permet de passer de
%   la langue courante à la langue spécifiée en argument.
% \end{DescribeMacro}
%
% \begin{DescribeEnv}{otherlanguage}
%   L'environnement |otherlanguage| de \pkg{babel} permet de faire
%   la même chose que la commande \cmd{\selectlanguage}, sauf que le
%   changement de langue est local à l'environnement --- utile pour
%   les brefs changements de langue.
% \end{DescribeEnv}
%
% Si vous n'êtes pas autrement familier avec le paquetage
% \pkg{babel}, consultez sa %
% \doc{babel}{http://texdoc.net/pkg/babel/}. %
% Celle-ci est éclatée en un
% document principal, pour le c{\oe}ur du paquetage et plusieurs
% autres pour les fonctionnalités propres à une langue: %
% \doc[anglais]{english}{http://texdoc.net/pkg/babel-english}, %
% \doc[français]{frenchb}{http://texdoc.net/pkg/babel-french}, %
% etc. Consultez au moins les documents consacrés aux
% langues utilisées dans votre thèse ou mémoire. Le plus simple
% consiste sans doute à consulter en ligne sur CTAN les
% \link{http://mirrors.ctan.org/tex-archive/macros/latex/required/babel/contrib/}{%
% documents spécifiques par langue}.
% \changes{4.4}{2017-05-14}{Correction d'une url vers la documentation
%   de babel.}
%
% \begin{DescribeMacro}{\nombre}
%   Le paquetage \pkg{numprint} étant chargé dans la classe avec
%   \pkg{babel}, vous pouvez utiliser la commande \cmd{\nombre} pour
%   formater automatiquement les nombres. Par exemple, le résultat de
%   |\nombre{123456789}| est \nombre{123456789}.
% \end{DescribeMacro}
%
% Les utilisateurs de {\XeLaTeX} qui souhaiteraient plutôt utiliser le
% plus récent paquetage \pkg{polyglossia} \citep{polyglossia} peuvent
% empêcher le chargement de \pkg{babel} avec l'option |nobabel| de la
% classe. Ils devront toutefois charger et configurer
% \pkg{polyglossia} eux-mêmes dans l'entête de leur document. Ce
% paquetage est moins évolué que \pkg{babel} pour la typographie
% française.
%
% \section{Bibliographie}
% \label{sec:bibliographie}
%
% \begin{DescribeMacro}{\bibliography}
%   Tel qu'expliqué au chapitre~8 de \citet{formation-latex-ul},
%   il est fortement recommandé d'utiliser {\BibTeX} pour la
%   préparation de la bibliographie d'un document. Celle-ci est
%   insérée dans le document à l'endroit où apparait la commande
%   \cmd{\bibliography} dans le code source. Cette commande prend en
%   arguments les noms des bases de données bibliographiques séparés
%   par des virgules.
% \end{DescribeMacro}
%
% \subsection{Mise en forme des citations}
% \label{sec:bibliographie:options}
%
% Dans la classe \class{ulthese}, la mise en forme des citations est
% confiée au paquetage \pkg{natbib} (à moins que l'option |nonatbib|
% ne soit spécifiée). Le paquetage est chargé avec les options par
% défaut, soit |round|, |semicolon| et |authoryear|. Pour spécifier
% d'autres options, vous avez deux possibilités:
% \begin{enumerate}
% \item utiliser l'option |nonatbib| de la classe et ensuite charger
%   explicitement \pkg{natbib} avec ses options;
% \item
%   \begin{DescribeMacro}{\setcitestyle}
%     utiliser la commande \cmd{\setcitestyle} pour passer de
%     nouvelles options à \pkg{natbib}.
%   \end{DescribeMacro}
% \end{enumerate}
%
% Par exemple, pour utiliser un style de citation numérique où le
% numéro de la référence se trouve entre crochets, vous pouvez
% procéder de l'une ou l'autre des deux manières suivantes:
% \begin{quote}
%   |\documentclass[nonatbib]{ulthese}| \\
%   |\usepackage[numbers,square]{natbib}| \\
%   |...| \\
%   |\bibliographystyle{plain-fr}|
% \end{quote}
% ou
% \begin{quote}
%   |\documentclass{ulthese}| \\
%   |...| \\
%   |\setcitestyle{numbers,square}| \\
%   |\bibliographystyle{plain-fr}|
% \end{quote}
%
% Consultez la %
% \doc{natbib}{http://texdoc.net/pkg/natbib/} %
% de \pkg{natbib} pour les détails.
%
% \subsection{Bibliographies multiples}
% \label{sec:bibliographie:multiples}
%
% La rédaction de la thèse ou du mémoire par articles est une pratique
% qui gagne en popularité. Elle consiste à remplacer le corps
% principal du manuscrit par un ou des articles scientifiques,
% habituellement à raison d'un par chapitre. Dans de tels cas, on
% souhaitera utiliser une bibliographie distincte pour chacun de ces
% chapitres. L'ajout de |bibchapitre| ou de |bibsection| dans les
% options de la classe permet de réaliser cette mise en forme
% particulière.
%
% La classe fait appel au paquetage \pkg{chapterbib}
% \citep{chapterbib} pour composer des bibliographies multiples. La
% \doc{chapterbib}{http://texdoc.net/pkg/chapterbib/} du paquetage
% explique la procédure à suivre pour obtenir une bibliographie par
% chapitre avec {\BibTeX}. En résumé:
% \begin{enumerate}
% \item chaque chapitre comportant sa propre bibliographie doit
%   absolument être inséré dans le document principal avec la commande
%   \cmd{\include} \citep[section~3.5]{formation-latex-ul};
% \item chaque fichier visé doit contenir une commande
%   \cmd{\bibliographystyle} (\autoref{sec:bibliographie:style}) et
%   une commande \cmd{\bibliography};
% \item on obtient les diverses bibliographies en compilant avec
%   {\BibTeX} les fichiers individuels, et non le document maître.
% \end{enumerate}
%
% La plupart des auteurs préféreront sans doute l'effet de l'option
% |bibsection|, où la bibliographie apparait comme une section normale
% à la fin d'un chapitre. Lorsque \pkg{babel} est chargé, la section
% sera numérotée. Pour supprimer la numérotation, insérer dans le
% préambule du document la commande
% \begin{quote}
%   |\addto\extras|\meta{langue}|{\renewcommand{\bibsection}{%| \\
%   |  \section*{\bibname}\prebibhook}}|
% \end{quote}
% où \meta{langue} est la langue par défaut du document
% (\autoref{sec:langues}).
%
% Le gabarit |gabarit-doctorat-articles| fournit la structure de base
% d'une thèse par articles.
%
% \subsection{Style de la bibliographie}
% \label{sec:bibliographie:style}
%
% \begin{DescribeMacro}{\bibliographystyle}
%   Le format général de la bibliographie est contrôlé par un
%   \emph{style} choisi avec la commande \cmd{\bibliographystyle} dans
%   le préambule du document ou dans les fichiers de chaque chapitre
%   dans une thèse ou un mémoire par articles. Les styles standards de
%   {\LaTeX} sont |plain|, |unsrt|, |alpha| et |abbrv|.
% \end{DescribeMacro}
%
% Le paquetage \pkg{natbib} chargé par défaut par la classe supporte
% le style de citation auteur-année fréquemment employé en sciences
% naturelles, plusieurs commandes de citation, un grand nombre de
% styles de bibliographie ainsi que des entrées spécifiques pour les
% numéros ISBN et les URL. Le paquetage fournit des styles de
% bibliographie |plainnat|, |unsrtnat| et |abbrvnat| similaires aux
% styles standards, mais plus complets. Il existe des
% \link{http://mirrors.ctan.org/biblio/bibtex/contrib/bib-fr/}{versions
% francisées} de ces styles (et de quelques autres) dans CTAN.
%
% Le paquetage \pkg{francais-bst} \citep{francais-bst} fournit une
% feuille de style compatible avec \pkg{natbib} permettant de composer
% des bibliographies auteur-année respectant les normes de typographie
% française proposées dans \cite{Malo:1996}. Pour utiliser ce style,
% spécifier dans le préambule du document LaTeX
% \begin{quote}
%   |\bibliographystyle{francais}|
% \end{quote}
%
% Autrement, la FESP n'a pas d'exigences particulières quant à la
% présentation de la bibliographie (présentation du titre, des auteurs
% et autres informations bibliographiques).
%
% \section{Police de caractères du document}
% \label{sec:police}
%
% Les documents {\LaTeX} sont facilement reconnaissables par leur
% police de caractères par défaut, {\fontfamily{cmr}\selectfont
% Computer Modern}. Avec toute distribution {\LaTeX} moderne, il est
% maintenant simple d'utiliser l'une ou l'autre des polices PostScript
% standards. D'ailleurs la classe \class{ulthese} utilise la police
% sans empattements \textsf{Helvetica} pour composer les pages de
% titre.
%
% La FESP permet l'utilisation des polices Times, Palatino (la police
% du présent document) et Lucida~Bright dans les thèse et mémoires.
% \citet[section~10.2]{formation-latex-ul} explique comment utiliser ces
% polices dans votre document.
%
% \section{Gabarits}
% \label{sec:gabarits}
%
% Les gabarits livrés avec la classe comportent des commentaires
% succincts pour vous guider dans la préparation de votre manuscrit.
% Les sections suivantes fournissent des détails additionnels, et ce,
% dans l'ordre où les commandes apparaissent dans les gabarits.
%
% \begin{rem}
%   Il n'y a pas de gabarit spécifique pour un examen de doctorat. On
%   utilise le gabarit de thèse approprié en ajoutant simplement
%   l'option |examen| dans la commande \cmd{\documentclass}.
% \end{rem}
%
% \subsection{Encodage des fichiers}
%
% Composer en {\LaTeX} de longs textes dans une langue ayant recours
% aux signes diacritiques devient rapidement pénible si l'on utilise
% des commandes telles que |\'e|, |\`a| ou |\"o| pour entrer des lettres
% accentuées. Afin de pouvoir plutôt entrer directement |é|, |à| ou
% |ö|, {\LaTeX} doit être configuré pour reconnaître les lettres
% accentuées. C'est le rôle du paquetage \pkg{inputenc}
% \citep{inputenc}.
%
% Il existe plusieurs manières différentes d'encoder --- ou
% d'enregistrer --- les lettres accentuées et autres caractères
% spéciaux (comme, par exemple, le symbole de l'euro) dans un
% ordinateur. La méthode la plus répandue et celle standard sur les
% versions récentes des systèmes d'exploitation Linux et macOS est
% l'UTF-8 de la norme
% \link{http://fr.wikipedia.org/wiki/Unicode}{Unicode}. Les gabarits
% sont livrés dans ce type d'encodage.
%
% La déclaration
% \begin{quote}
%   |\usepackage[utf8]{inputenc}|
% \end{quote}
% dans le préambule assure que {\LaTeX} traitera correctement des
% fichiers source encodés en UTF-8.
%
% La norme Unicode n'est pas aussi uniformément supportée par Windows.
% Selon l'éditeur de texte employé et la version du système
% d'exploitation, il peut être nécessaire d'utiliser les normes
% d'encodage %
% \link{http://fr.wikipedia.org/wiki/ISO_8859-1}{ISO~8859-1} %
% (ou Latin-1; option |latin1| de \pkg{inputenc}), %
% \link{http://fr.wikipedia.org/wiki/ISO_8859-15}{ISO~8859-15} %
% (ou Latin-9; option |latin9|) ou %
% \link{http://fr.wikipedia.org/wiki/Windows-1252}{Windows-1252} %
% (options |cp1252| ou |ansinew|).
%
% La situation est plus simple avec {\XeLaTeX} puisqu'il gère
% nativement Unicode. Le paquetage \pkg{inputenc} est non seulement
% inutile, mais incompatible avec {\XeLaTeX}. C'est pourquoi, dans les
% gabarits, \pkg{inputenc} est chargé seulement lorsque {\XeTeX}
% n'est pas le moteur employé pour compiler le document.
%
% \subsection{Paquetages additionnels}
%
% Tel qu'expliqué à la \autoref{sec:paquetages}, la classe
% charge déjà quelques paquetages. Cependant, il est fort probable que
% vous devrez en charger d'autres pour composer votre document. Les
% gabarits prévoient un endroit pour le chargement de paquetages
% additionnels. Il est recommandé d'inscrire vos commandes
% \cmd{\usepackage} à cet endroit afin de respecter un certain ordre de
% chargement; voir ci-dessous.
%
% Si vous utilisez un paquetage non standard dans les distributions
% courantes (\TeX~Live, Mac\TeX, MiK\TeX), vous devez le fournir avec
% le code source de votre document lors du dépôt final.
%
% \subsection{Changement de police de caractères}
%
% Les gabarits comportent des déclarations types pour utiliser les
% polices Palatino ou Times sous {\LaTeX} ou, sous {\XeLaTeX}, leurs
% équivalents Pagella et Termes du projet
% \link{http://www.gust.org.pl/projects/e-foundry/tex-gyre/}{TeX~Gyre}.
%
% \subsection{Hyperliens}
% \label{sec:hyperref}
%
% Le paquetage \pkg{hyperref} \citep{hyperref} permet de transformer
% toutes les références en hyperliens cliquables lorsque le document
% est produit avec pdf{\LaTeX} ou {\XeLaTeX}. L'interaction de ce
% paquetage avec les autres est parfois (voire souvent) délicate. Pour
% cette raison, il est habituellement nécessaire de charger
% \pkg{hyperref} en tout dernier. C'est pourquoi il n'est pas chargé
% dans la classe, mais plutôt dans les gabarits. Prenez soin de
% maintenir le dernier rang de chargement lors de l'édition d'un
% gabarit.
%
% \begin{DescribeMacro}{\hyperrefsetup}
%   La configuration du paquetage dans les gabarits fait en sorte que
%   les liens sont simplement signalés par une couleur de texte
%   légèrement contrastante. L'utilisation de couleurs dans un
%   document requiert le paquetage |xcolor|, chargé par la classe. La
%   couleur de lien par défaut, |ULlinkcolor|, est définie dans la
%   classe; voir la \autoref{sec:couleurs}.
% \end{DescribeMacro}
%
% \subsection{Options de \pkg{babel}}
%
% \begin{DescribeMacro}{\frenchbsetup}
%   La commande \cmd{\frenchbsetup} de \pkg{babel} permet de contrôler
%   certains ajustements typographiques apportés par le paquetage en
%   mode français. Consultez la documentation de \pkg{babel} pour la
%   liste des options de configuration disponibles.
% \end{DescribeMacro}
%
% Les concepteurs de la classe \class{ulthese}  proposent trois
% ajustements dans les gabarits:
% \begin{enumerate}
% \item l'option |StandardItemizeEnv=true| évite que le mode français de
%   \pkg{babel} ne diminue l'espacement vertical dans les listes;
% \item l'option |ThinSpaceInFrenchNumbers=true| fait en sorte qu'une espace fine
%   sera utilisée comme séparateur des milliers dans les nombres plutôt
%   qu'une espace pleine;
% \item les options |og=«| et |fg=»| déclarent que les caractères « et
%   » utilisés dans le code source représentent les guillemets ouvrant
%   et fermant, respectivement. Cela évite de devoir utiliser les
%   commandes \cmd{\og} et \cmd{\fg} de \pkg{babel} tout en bénéficiant de
%   l'ajutement automatique des espaces autour des symboles.
% \end{enumerate}
% Vous devez évidemment désactiver ces ajustements si l'option
% |nobabel| est spécifiée au chargement de la classe.
%
% \subsection{Style de la bibliographie}
% \label{sec:bibliographystyle}
%
% Tel qu'expliqué à la \autoref{sec:bibliographie:style}, la commande
% \cmd{\bibliographystyle} permet de définir le mode de citation dans
% le texte et la présentation des notices bibliographiques. Cette
% commande devrait figurer dans les fichiers des chapitres
% individuels dans une thèse ou un mémoire par articles.
%
% \subsection{Déclarations des pages de titre}
%
% Les gabarits comportent toutes les déclarations nécessaires pour
% composer les pages de titre des divers types de thèse ou de mémoires.
% Vous devez remplacer les éléments se trouvant entre crochets <~> en
% respectant la forme indiquée. Assurez-vous de supprimer les
% caractères < et > afin qu'ils n'apparaissent pas sur les pages de titre de
% votre document.
%
% \subsection{Pages liminaires}
%
% \begin{DescribeMacro}{\frontmatter}
%   La commande \cmd{\frontmatter} déclare que {\LaTeX} doit considérer le
%   matériel qui suit comme des pages liminaires. En pratique, cela
%   résulte essentiellement en une numérotation des pages en chiffres
%   romains.
% \end{DescribeMacro}
%
% Les normes de présentation de la FESP édictent que les thèses et
% mémoires devraient comporter les pages liminaires suivantes, dans
% l'ordre:
% \begin{enumerate}
% \item la page de titre (obligatoire);
% \item la page frontispice (obligatoire);
% \item un résumé en français (obligatoire);
% \item un résumé en anglais (recommandé, mais non obligatoire);
% \item une table des matières (obligatoire);
% \item une liste des tableaux;
% \item une liste des figures;
% \item une liste des abbréviations et des sigles;
% \item une dédicace;
% \item une épigraphe;
% \item des remerciements;
% \item un avant-propos (obligatoire dans le cas d'une thèse ou d'un
%   mémoire par articles).
% \end{enumerate}
%
% \begin{rem}
%   Pour un examen de doctorat, on ne considère que la page de titre
%   comme obligatoire. L'option de classe |examen| supprime d'ailleurs
%   la page frontispice. Il est laissé aux auteurs le soin de
%   supprimer les autres pages liminaires des gabarits.
% \end{rem}
%
% Les commandes
% \begin{quote}
%   \cmd{\pagestitre} \\
%   \cmd{\tableofcontents} \\
%   \cmd{\listoftables} \\
%   \cmd{\listoffigures} \\
%   \cmd{\dedicace}\marg{texte} \\
%   \cmd{\epigraphe}\marg{texte}\marg{auteur}
% \end{quote}
% permettent de générer les pages correspondantes. Seules les deux
% dernières commandes admettent des arguments.
%
% \begin{DescribeMacro}{\chapter*}
%   Les résumés, la liste des abbréviations et des sigles, les
%   remerciements et l'avant-propos sont composés comme des chapitres
%   normaux, mais sans être numérotés. Il faut donc définir ces éléments
%   avec la commande \cmd{\chapter*}.
% \end{DescribeMacro}
%
% \begin{DescribeMacro}{\phantomsection}
%   \begin{DescribeMacro}{\addcontentsline}
%     Les sections declarées avec la commande \cmd{\chapter*}
%     n'apparaissent pas dans la table des matières. Comme les normes
%     de présentation de la FESP exigent que toutes les pages
%     liminaires y figurent, on fait suivre les commandes
%     \cmd{\chapter*}\marg{Titre} des commandes
%   \end{DescribeMacro}
% \end{DescribeMacro}
% \begin{quote}
%   |\phantomsection\addcontentsline{toc}{chapter}|\marg{Titre}
% \end{quote}
% Celles-ci ajoutent à la table des matières (|toc|) une section de
% niveau |chapter| dont le titre est \meta{Titre}. La commande
% \cmd{\phantomsection} est rendue nécessaire (ou recommandée) par le
% paquetage \pkg{hyperref}.
%
% \subsection{Corps du document}
%
% \begin{DescribeMacro}{\mainmatter}
%   La commande \cmd{\mainmatter} délimite le début du corps du document. La
%   numérotation des pages passe en chiffres arabes.
% \end{DescribeMacro}
%
% Le corps du document devrait normalement compter une introduction (non
% numérotée), un développement divisé en chapitres (numérotés) et une
% conclusion (non numérotée).
%
% \subsection{Annexes}
%
% \begin{DescribeMacro}{\appendix}
%   Si la thèse ou le mémoire comporte une ou plusieurs annexes,
%   composer celles-ci comme des chapitres normaux insérés dans le
%   document maître après la commande \cmd{\appendix}. Cette commande a
%   pour effet de passer d'un mode de numération numérique (1, 1.1, 2, 2.1,
%   \dots) à un mode alphanumérique (A, A.1, B, B.1, \dots).
% \end{DescribeMacro}
%
% \subsection{Bibliographie}
%
% La bibliographie figure normalement à la toute fin du document, sous
% forme de chapitre non numéroté. Tel qu'expliqué à la
% \autoref{sec:bibliographie:multiples}, la commande
% \cmd{\bibliography} doit plutôt se trouver dans les fichiers des
% chapitres dans une thèse ou un mémoire par articles.
%
% \section{Aide additionnelle}
%
% Pour obtenir de l'aide additionnelle sur l'utilisation de la classe
% \class{ulthese} (et non sur celle de {\LaTeX} en général), prière
% de consulter d'abord
% \begin{enumerate}
% \item le \link{http://www.theses.ulaval.ca/wiki/}{WikiThèse} de
%   l'Université Laval, en particulier la
%   \link{http://www.theses.ulaval.ca/wiki/index.php?title=FAQ}{Foire aux questions};
% \item les
%   \link{http://listes.ulaval.ca/listserv/archives/ulthese-aide.html}{archives}
%   de la liste de distribution \texttt{ulthese-aide}.
% \end{enumerate}
% Si la réponse à votre question ne se trouve ni dans le wiki, ni dans
% les archives, alors écrire à l'adresse
% \href{mailto:ulthese-aide@listes.ulaval.ca}{\url{ulthese-aide@listes.ulaval.ca}}.
%
% \StopEventually{
%   \begin{thebibliography}{15}
%     \addcontentsline{toc}{section}{Références}
%   \bibitem[{Arseneau(2010)}]{chapterbib}
%     Arseneau, D. 2010,
%     {\selectlanguage{english}\emph{chapterbib. Multiple
%     bibliographies in {\LaTeX}}}.
%     URL~\url{http://www.ctan.org/pkg/chapterbib/}.
%
%   \bibitem[{Braams et Bezos(2016)}]{babel}
%     Braams, J. et J.~Bezos. 2016,
%     {\selectlanguage{english}\emph{Babel}}.
%     URL~\url{http://www.ctan.org/pkg/babel/}.
%
%   \bibitem[{Carlisle et {The \LaTeX3\ Project}(2016)}]{graphicx}
%     Carlisle, D. et {The \LaTeX3\ Project}. 2016,
%     {\selectlanguage{english}\emph{Packages in the `graphics'
%     Bundle}}. URL~\url{http://www.ctan.org/pkg/graphics/}.
%
%   \bibitem[{Charette(2015)}]{polyglossia}
%     Charette, F. 2015, {\selectlanguage{english}\emph{Polyglossia:
%     An Alternative to Babel for {\XeLaTeX} and {\LuaLaTeX}}}.
%     URL~\url{http://www.ctan.org/pkg/polyglossia/}, current
%     maintainer Arthur Reutenauer.
%
%   \bibitem[{Daly(2010)}]{natbib}
%     Daly, P.~W. 2010, {\selectlanguage{english}\emph{Natural
%     Sciences Citations and References}}.
%     URL~\url{http://www.ctan.org/pkg/natbib/}.
%
%   \bibitem[{Goulet(2013)}]{francais-bst}
%     Goulet, V. 2013, {\selectlanguage{french}«Paquetage
%     \pkg{francais-bst}»},
%     URL~\url{http://www.ctan.org/pkg/francais-bst/}.
%
%   \bibitem[{Goulet(2016)}]{formation-latex-ul}
%     Goulet, V. 2016, {\selectlanguage{french}\emph{Rédaction avec
%     {\LaTeX}}}, document libre sous contrat Creative Commons.
%     ISBN~978-2-9811416-7-5.
%     URL~\url{https://ctan.org/pkg/formation-latex-ul}.
%
%   \bibitem[{Harders(2012)}]{numprint}
%     Harders, H. 2012, {\selectlanguage{english}\emph{The
%     \pkg{numprint} package}}.
%     URL~\url{http://www.ctan.org/pkg/numprint/}.
%
%   \bibitem[{Jeffrey et Mittelbach(2015)}]{inputenc}
%     Jeffrey, A. et F.~Mittelbach. 2015,
%     {\selectlanguage{english}\emph{inputenc.sty}}.
%     URL~\url{http://www.ctan.org/pkg/inputenc/}.
%
%   \bibitem[{Kern(2016)}]{xcolor}
%     Kern, D.~U. 2016, {\selectlanguage{english}\emph{Extending
%     {\LaTeX}’s color facilities: the \pkg{xcolor} package}}.
%     URL~\url{http://www.ctan.org/pkg/xcolor/}.
%
%   \bibitem[{Malo(1996)}]{Malo:1996}
%     Malo, M. 1996, {\selectlanguage{french}\emph{Guide de la
%     communication écrite au cégep, à l'université et en
%     entreprise}}, Québec Amérique. ISBN~978-2-8903-7875-9.
%
%   \bibitem[{Rahtz et Oberdiek(2017)}]{hyperref}
%     Rahtz, S. et H.~Oberdiek. 2017,
%     {\selectlanguage{english}\emph{Hypertext marks in {\LaTeX}: a
%     manual for \pkg{hyperref}}}.
%     URL~\url{http://www.ctan.org/pkg/hyperref/}.
%
%   \bibitem[{Robertson et Hosny(2017)}]{fontspec}
%     Robertson, W. et K.~Hosny. 2017,
%     {\selectlanguage{english}\emph{The \pkg{fontspec} package: Font
%     selection for {\XeLaTeX} and {\LuaLaTeX}}}.
%     URL~\url{http://www.ctan.org/pkg/fontspec/}.
%
%   \bibitem[{Robertson et coll.(2017)Robertson, Stephani et
%     Hosny}]{unicode-math}
%     Robertson, W., P.~Stephani et K.~Hosny. 2017,
%     {\selectlanguage{english}\emph{Experimental {U}nicode
%     Mathematical Typesetting: The \pkg{unicode-math} Package}}.
%     ULR~\url{http://www.ctan.org/pkg/unicode-math/}.
%
%   \bibitem[{Wilson(2016)}]{memoir}
%     Wilson, P. 2016, {\selectlanguage{english}\emph{The Memoir Class
%     for Configurable Typesetting}}, 8{\ieme} éd., The Herries Press.
%     URL~\url{http://www.ctan.org/pkg/memoir/}, maintained by Lars
%     Madsen.
%   \end{thebibliography}
%   \PrintChanges
% }
%
% ^^A Début du code de la classe
%
% \appendix
% \section{Mise en {\oe}uvre}
% \label{sec:meo}
%
% Cette annexe passe en revue le code {\TeX} et {\LaTeX} de la
% classe. Elle ne risque d'intéresser que les personnes qui souhaitent
% explorer comment la classe est programmée.
%
% \subsection{Tests et valeurs booléennes}
% \changes{4.4}{2017-06-01}{Tests et valeurs booléennes réalisés sans ifthen}
%
% Le paquetage \pkg{ifxetex} permet de tester si un document est
% compilé avec \XeTeX.
%    \begin{macrocode}
%<*class>
\RequirePackage{ifxetex}
%    \end{macrocode}
% Nous définissons ici toutes les valeurs booléennes requises par la
% classe.
%    \begin{macrocode}
\newif\ifUL@babel       \UL@babeltrue        % charger babel?
\newif\ifUL@natbib      \UL@natbibtrue       % charger natbib?
\newif\ifUL@chapterbib  \UL@chapterbibfalse  % charger chapterbib?
\newif\ifUL@sectionbib  \UL@sectionbibfalse  % option sectionbib de chapterbib?
\newif\ifUL@isthesis                         % programme est une thèse?
\newif\ifUL@iscotutelle \UL@iscotutellefalse % thèse en cotutelle?
\newif\ifUL@isexam      \UL@isexamfalse      % examen de doctorat?
\newif\ifUL@hassubtitle \UL@hassubtitlefalse % document a un sous-titre?
%    \end{macrocode}
%
% \subsection{Options de la classe}
%
% Il y a cinq grandes catégories d'options propres à la classe: la
% possibilité d'empêcher le chargement du paquetage \pkg{natbib};
% la possibilité d'empêcher le chargement du paquetage \pkg{babel};
% la taille de la police de caractères en points; le type de grade; la
% déclaration qu'il s'agit d'un examen de doctorat.
%
% \begin{macro}{nonatbib}
%   L'option |nonatbib| permet d'empêcher la classe de charger le
%   paquetage \pkg{natbib} en cas d'incompatibilité avec d'autres
%   paquetages spécialisés de mise en forme de la bibliographie.
%    \begin{macrocode}
\DeclareOption{nonatbib}{\UL@natbibfalse}
%    \end{macrocode}
% \end{macro}
%
% \begin{macro}{bibchapitre}
%   \begin{macro}{bibsection}
%     L'option |bibchapitre| entraine le chargement du paquetage
%     \pkg{chapterbib} qui permet de composer des bibliographies
%     multiples dans un document, habituellement une par
%     chapitre. Avec l'option |bibsection|, la bibliographie se
%     présente sous la forme d'une section plutôt que d'un chapitre.
%     La seconde option implique la première.
%    \begin{macrocode}
\DeclareOption{bibchapitre}{\UL@chapterbibtrue}
\DeclareOption{bibsection}{\UL@chapterbibtrue\UL@sectionbibtrue}
%    \end{macrocode}
%   \end{macro}
% \end{macro}
%
% \begin{macro}{nobabel}
%   L'option |nobabel| permet d'empêcher la classe de charger le
%   paquetage \pkg{babel}. Cette option peut s'avérer utile pour
%   les utilisateurs de {\XeLaTeX} qui souhaitent plutôt utiliser
%   \pkg{poyglossia} pour le traitement des langues dans leur
%   document.
%    \begin{macrocode}
\DeclareOption{nobabel}{\UL@babelfalse}
%    \end{macrocode}
% \end{macro}
%
% \begin{macro}{10pt,11pt,12pt}
%   Les valeurs possibles pour la taille de la police de caractères
%   sont |10pt|, |11pt| et |12pt|. Cette option est gérée au niveau de
%   la classe afin de s'assurer que les divers éléments sur les pages
%   de titre sont toujours de la même taille. La taille de la police
%   par défaut permet de déterminer si, par exemple, le titre du
%   document doit être dans la taille \cmd{\Huge}, \cmd{\huge} ou \cmd{\LARGE} de
%   \class{memoir}.
%
%   La taille de la police est passée à \class{memoir} et la macro
%   \cmd{\UL@ptsize} stocke la taille des caractères pour usage futur.
%    \begin{macrocode}
\newcommand*{\UL@ptsize}{}
\DeclareOption{10pt}{%
  \PassOptionsToClass{10pt}{memoir}
  \renewcommand*{\UL@ptsize}{10}}
\DeclareOption{11pt}{%
  \PassOptionsToClass{11pt}{memoir}
  \renewcommand*{\UL@ptsize}{11}}
\DeclareOption{12pt}{%
  \PassOptionsToClass{12pt}{memoir}
  \renewcommand*{\UL@ptsize}{12}}
%    \end{macrocode}
% \end{macro}
%
% \begin{macro}{PhD,MSc,MA,...}
%   Définition du type de grade et si la thèse ou le mémoire est
%   multifacultaire, effectué en cotutelle, en bidiplomation ou en
%   extension. Certaines options ne sont valides que pour une thèse ou
%   que pour un mémoire. La page de titre des programmes en extension
%   comporte une mention «offert en extension» ou «offerte en
%   extension» selon qu'il s'agit d'un doctorat ou d'une maîtrise; le
%   bon terme est défini avec l'option correspondante.
%    \begin{macrocode}
\newcommand*{\UL@typenum}{}
\DeclareOption{LLD}{%
  \UL@isthesistrue
  \renewcommand*{\UL@typenum}{0}
  \newcommand*{\UL@degree}{Docteur en droit (LL.~D.)}}
\DeclareOption{DMus}{%
  \UL@isthesistrue
  \renewcommand*{\UL@typenum}{0}
  \newcommand*{\UL@degree}{Docteur en musique (D.~Mus.)}}
\DeclareOption{DPsy}{%
  \UL@isthesistrue
  \renewcommand*{\UL@typenum}{0}
  \newcommand*{\UL@degree}{Docteur en psychologie (D.~Psy.)}}
\DeclareOption{DThP}{%
  \UL@isthesistrue
  \renewcommand*{\UL@typenum}{0}
  \newcommand*{\UL@degree}{Docteur en th\'eologie pratique (D.~Th.~P.)}}
\DeclareOption{PhD}{%
  \UL@isthesistrue
  \renewcommand*{\UL@typenum}{0}
  \newcommand*{\UL@degree}{Philosophi{\ae} doctor (Ph.~D.)}}
\DeclareOption{MATDR}{%
  \UL@isthesisfalse
  \renewcommand*{\UL@typenum}{0}
  \newcommand*{\UL@degree}{Ma\^itre en am\'enagement du territoire %
    et d\'eveloppement r\'egional (M.ATDR)}}
\DeclareOption{MArch}{%
  \UL@isthesisfalse
  \renewcommand*{\UL@typenum}{0}
  \newcommand*{\UL@degree}{Ma\^itre en architecture (M.~Arch.)}}
\DeclareOption{MA}{%
  \UL@isthesisfalse
  \renewcommand*{\UL@typenum}{0}
  \newcommand*{\UL@degree}{Ma\^itre \`es arts (M.A.)}}
\DeclareOption{LLM}{%
  \UL@isthesisfalse
  \renewcommand*{\UL@typenum}{0}
  \newcommand*{\UL@degree}{Ma\^itre en droit (LL.~M.)}}
\DeclareOption{MErg}{%
  \UL@isthesisfalse
  \renewcommand*{\UL@typenum}{0}
  \newcommand*{\UL@degree}{Ma\^itre en ergoth\'erapie (M.~Erg.)}}
\DeclareOption{MMus}{%
  \UL@isthesisfalse
  \renewcommand*{\UL@typenum}{0}
  \newcommand*{\UL@degree}{Ma\^itre en musique (M.~Mus.)}}
\DeclareOption{MPht}{%
  \UL@isthesisfalse
  \renewcommand*{\UL@typenum}{0}
  \newcommand*{\UL@degree}{Ma\^itre en physioth\'erapie (M.~Pht.)}}
\DeclareOption{MSc}{%
  \UL@isthesisfalse
  \renewcommand*{\UL@typenum}{0}
  \newcommand*{\UL@degree}{Ma\^itre \`es sciences (M.~Sc.)}}
\DeclareOption{MScGeogr}{%
  \UL@isthesisfalse
  \renewcommand*{\UL@typenum}{0}
  \newcommand*{\UL@degree}{Ma\^itre en sciences g\'eographiques (M.~Sc.~g\'eogr.)}}
\DeclareOption{MServSoc}{%
  \UL@isthesisfalse
  \renewcommand*{\UL@typenum}{0}
  \newcommand*{\UL@degree}{Ma\^itre en service social (M.~Serv.~soc.)}}
\DeclareOption{MPsEd}{%
  \UL@isthesisfalse
  \renewcommand*{\UL@typenum}{0}
  \newcommand*{\UL@degree}{Ma\^itre en psycho\'education (M.~Ps.~\'ed.)}}
\DeclareOption{multifacultaire}{%
  \ifUL@isthesis
    \renewcommand*{\UL@typenum}{1}
  \else
    \ClassError{ulthese}{%
      Incompatible option multifacultaire}
    {Use this option with a doctorate degree only.}
  \fi}
\DeclareOption{cotutelle}{%
  \ifUL@isthesis
    \renewcommand*{\UL@typenum}{2}
    \UL@iscotutelletrue
  \else
    \ClassError{ulthese}{%
      Incompatible option cotutelle}
    {Use this option with a doctorate degree only.}
  \fi}
\DeclareOption{bidiplomation}{%
  \ifUL@isthesis
    \ClassError{ulthese}{%
      Incompatible option bidiplomation}
    {Use this option with a master degree only.}
  \else
    \renewcommand*{\UL@typenum}{2}
  \fi}
\DeclareOption{extensionUdeS}{%
  \ifUL@isthesis
    \renewcommand*{\UL@typenum}{3}
    \newcommand*{\UL@offered}{offert}
    \newcommand*{\UL@extensionat}{Universit\'e de Sherbrooke}
    \newcommand*{\UL@extensionloc}{Sherbrooke, Canada}
  \else
    \ClassError{ulthese}{%
      Incompatible option extensionUdeS}
    {Use this option with a doctorate degree only.}
  \fi}
\DeclareOption{extensionUQO}{%
  \ifUL@isthesis
    \renewcommand*{\UL@typenum}{3}
    \newcommand*{\UL@offered}{offert}
    \newcommand*{\UL@extensionat}{Universit\'e du Qu\'ebec en Outaouais}
    \newcommand*{\UL@extensionloc}{Gatineau, Canada}
  \else
    \ClassError{ulthese}{%
      Incompatible option extensionUQO}
    {Use this option with a doctorate degree only.}
  \fi}
\DeclareOption{extensionUQAC}{%
  \ifUL@isthesis
    \ClassError{ulthese}{%
      Incompatible option extensionUQAC}
    {Use this option with a master degree only.}
  \else
    \renewcommand*{\UL@typenum}{3}
    \newcommand*{\UL@offered}{offerte}
    \newcommand*{\UL@extensionat}{Universit\'e du Qu\'ebec \`a Chicoutimi}
    \newcommand*{\UL@extensionloc}{Chicoutimi, Canada}
  \fi}
%    \end{macrocode}
% \end{macro}
%
% \begin{macro}{examen}
%   L'option |examen| change l'appellation «Thèse» sur la couverture
%   pour «Examen de doctorat» et supprime la page frontispice. Elle
%   n'est compatible qu'avec l'une des options de thèse, autrement un
%   message d'erreur est émis.
%    \begin{macrocode}
\DeclareOption{examen}{%
  \ifUL@isthesis
    \UL@isexamtrue
  \else
    \ClassError{ulthese}{%
      Incompatible option examen}
    {Use this option with a thesis type only.}
  \fi}
%    \end{macrocode}
% \end{macro}
%
% \subsection{Chargement de la classe \class{memoir}}
%
% Toutes les options de la classe sont passées à \class{memoir}. Le
% format de papier et la taille de police par défaut sont, dans
% l'ordre, |letterpaper| et |11pt|. On vérifie qu'un type de grade a
% bien été déclaré. L'option de \class{memoir} |oneside| est
% explicitement déclarée afin d'éviter toute tentative de passer outre
% à cette exigence de la FESP.
%    \begin{macrocode}
\DeclareOption*{\PassOptionsToClass{\CurrentOption}{memoir}}
\ExecuteOptions{11pt,letterpaper}
\ProcessOptions\relax
\ifx\UL@typenum\empty
  \ClassError{ulthese}{%
    No thesis type specified}
    {Declare the thesis type as a class option.}
\fi
\LoadClass[oneside]{memoir}
%    \end{macrocode}
%
% \subsection{Paquetages requis}
%
% La classe s'efforce de charger un minimum de paquetages afin d'éviter
% les conflits potentiels.
%
% {\XeLaTeX} requiert le paquetage \pkg{fontspec} pour le
% traitement des polices. Le paquetage \pkg{unicode-math} facilite
% également le traitement des polices et des symboles mathématiques
% avec ce moteur. Sous {\LaTeX}, il est aujourd'hui préférable
% d'utiliser les polices T1.
%    \begin{macrocode}
\ifxetex
  \RequirePackage{fontspec}
  \RequirePackage{unicode-math}
  \defaultfontfeatures{Ligatures=TeX}
\else
  \RequirePackage[T1]{fontenc}
\fi
%    \end{macrocode}
%
% Les paquetages \pkg{natbib} et \pkg{chapterbib} doivent être chargés
% avant \pkg{babel} pour bien fonctionner, le cas échéant. Tel que
% précisé dans la documentation de \pkg{natbib}, l'option |sectionbib|
% est passée à ce paquetage lorsqu'il est chargé, ou à
% \pkg{chapterbib} autrement.
%    \begin{macrocode}
\ifUL@natbib
  \ifUL@sectionbib
    \PassOptionsToPackage{sectionbib}{natbib}
  \fi
  \RequirePackage[round,semicolon,authoryear]{natbib}
\fi
\ifUL@chapterbib
  \ifUL@sectionbib
     \ifUL@natbib\else
       \PassOptionsToPackage{sectionbib}{chapterbib}
     \fi
  \fi
  \RequirePackage{chapterbib}
\fi
%    \end{macrocode}
%
% Le support pour les langues autres que l'anglais est offert par le
% paquetage \pkg{babel} --- à moins que l'option |nobabel| n'ait
% été spécifiée au chargement de la classe. Les langues sont passées
% en option de la classe, et non du paquetage. Le paquetage
% \pkg{numprint} est requis par \pkg{babel} pour la définition
% de la commande de mise en forme des nombres \cmd{\nombre}.
%    \begin{macrocode}
\ifUL@babel
  \RequirePackage{babel}
  \RequirePackage[autolanguage]{numprint}
\fi
%    \end{macrocode}
%
% L'insertion du logo de l'Université sur la page de titre requiert
% \pkg{graphicx}. Les coloration des hyperliens requiert \pkg{xcolor}.
% \autoref{sec:hyperref}).
%    \begin{macrocode}
\RequirePackage{graphicx}
\RequirePackage{xcolor}
%    \end{macrocode}
%
% La commande \cmd{\textcopyright} utilisée sur la page de titre requiert le
% paquetage \pkg{textcomp} pour obtenir un beau signe de copyright.
%    \begin{macrocode}
\RequirePackage{textcomp}
%    \end{macrocode}
%
% \subsection{Paquetage incompatible}
%
% Le chargement du paquetage \pkg{geometry} avec la classe
% \class{memoir} modifie les marges du document. Pour cette raison,
% \pkg{geometry} est déclaré incompatible avec la classe.
%    \begin{macrocode}
\AtBeginDocument{%
  \@ifpackageloaded{geometry}{%
    \ClassError{ulthese}{%
      Package geometry is incompatible with this class}
    {Use the memoir class facilities to change the page layout.}}{\relax}}
%    \end{macrocode}
%
% \subsection{Couleur des hyperliens}
% \label{sec:couleurs}
%
% Le paquetage \pkg{hyperref} est chargé dans les gabarits afin de
% demeurer le dernier paquetage chargé; voir la
% \autoref{sec:hyperref}. La classe définit néanmoins une couleur standard pour
% les hyperliens, une teinte de bleu assez foncée pour être à fois
% visible en couleur et peu contrastante si le document est imprimé en
% noir et blanc.
%    \begin{macrocode}
\definecolor{ULlinkcolor}{rgb}{0,0,0.3}
%    \end{macrocode}
%
% \subsection{Marges}
%
% Les marges exigées par les normes de présentation de la FESP sont de
% 30~mm pour les marges gauche et droite et 25~mm pour les marges
% supérieure et inférieure. Le pied de page est placé de sorte que le
% folio de page se retrouve à 10~mm du bas de la page.
%    \begin{macrocode}
\setlrmarginsandblock{30mm}{30mm}{*}
\setulmarginsandblock{25mm}{25mm}{*}
\checkandfixthelayout[nearest]
\setlength{\footskip}{\lowermargin}
\addtolength{\footskip}{-10mm}
%    \end{macrocode}
%
% Comme les thèses et mémoires comportent normalement plusieurs pages
% liminaires, il arrive que des folios (en chiffres romains) dépassent
% dans la marge de droite dans la table des matières. Pour régler ce
% problème, nous augmentont la largeur de la boîte prévue pour les
% imprimer.
%    \begin{macrocode}
\renewcommand{\@pnumwidth}{3em}
\renewcommand{\@tocrmarg}{4em}
%    \end{macrocode}
%
% \subsection{Interligne}
%
% L'espacement entre les lignes est d'un interligne et demi.
% L'espacement «double» entre les paragraphes est fixé à
% |0.5\baselineskip| afin d'en arriver à une disposition agréable à
% l'{\oe}il. Le retrait de première ligne est supprimé puisque plus
% nécessaire suite à l'ajout de l'espacement entre les paragraphes.
%    \begin{macrocode}
\OnehalfSpacing
\setlength{\parskip}{0.5\baselineskip}
\setlength{\parindent}{0em}
%    \end{macrocode}
%
% La table des matières, la liste des tableaux et la liste des figures
% sont composées à interligne simple.
%    \begin{macrocode}
\renewcommand{\tocheadstart}{\SingleSpacing\chapterheadstart}
\renewcommand{\lotheadstart}{\SingleSpacing\chapterheadstart}
\renewcommand{\lofheadstart}{\SingleSpacing\chapterheadstart}
%    \end{macrocode}
%
% \subsection{Entêtes et pieds de page}
%
% Les règles pour les entêtes et pieds de page sont uniformes pour
% tout le document: aucun entête et folio au centre du pied de page.
% Ceci correspond au style standard |plain|.
%    \begin{macrocode}
\pagestyle{plain}
%    \end{macrocode}
%
% \subsection{Pages de titre}
%
% Le code pour traiter et composer la page de titre et la page
% frontispice constitue l'essentiel de la classe.
%
% \subsubsection{Famille et style de la police de caractères}
%
% Les pages de titre sont composées avec la police \textsf{Helvetica}
% (famille \texttt{phv} dans la classification NFSS) dans les
% tailles\footnote{%
%   La police Helvetica produite par {\LaTeX} est plus grande que
%   celle utilisée par Microsoft Word. Pour cette raison, les tailles
%   utilisées dans la classe sont toutes quelques points inférieures à
%   celles des gabarits Word.}%
% et les graisses présentées au \autoref{tab:polices}. La
% déclaration |\fontencoding{T1}| est nécessaire avec {\XeLaTeX} pour
% explicitement charger la même police que sous {\LaTeX}.
%
% \begin{table}
%   \centering
%   \begin{tabular}{ll}
%     \toprule
%     Élément          & Police \\
%     \midrule
%     Titre            & 17~points gras \\
%     Sous-titre       & 14~points gras \\
%     Auteur           & 12~points gras \\
%     Nom du programme & 12~points gras \\
%     Autres éléments  & 12~points normal \\
%     \bottomrule
%   \end{tabular}
%   \caption{Tailles et graisses de la police Helvetica des éléments
%     de la page de titre}
%   \label{tab:polices}
% \end{table}
%
%    \begin{macrocode}
\newcommand*{\UL@phvfamily}{\fontencoding{T1}\fontfamily{phv}\selectfont}
%    \end{macrocode}
% Les commandes sélectionnant ces polices sont adaptées selon la
% taille de police choisie pour le document afin d'être toujours
% identiques. Nous utilisons les déclarations de taille de police de
% la classe \class{memoir}, présentées au tableau~3.9 de sa
% documentation.
%    \begin{macrocode}
\ifnum\UL@ptsize=10\relax
  \newcommand*{\UL@fonttitle}{\normalfont\huge\bfseries\UL@phvfamily}
  \newcommand*{\UL@fontsubtitle}{\normalfont\LARGE\bfseries\UL@phvfamily}
  \newcommand*{\UL@fontauthor}{\normalfont\Large\bfseries\UL@phvfamily}
  \newcommand*{\UL@fontprogram}{\UL@fontauthor}
  \newcommand*{\UL@fontbase}{\normalfont\Large\UL@phvfamily}
\fi
\ifnum\UL@ptsize=11\relax
  \newcommand*{\UL@fonttitle}{\normalfont\LARGE\bfseries\UL@phvfamily}
  \newcommand*{\UL@fontsubtitle}{\normalfont\Large\bfseries\UL@phvfamily}
  \newcommand*{\UL@fontauthor}{\normalfont\large\bfseries\UL@phvfamily}
  \newcommand*{\UL@fontprogram}{\UL@fontauthor}
  \newcommand*{\UL@fontbase}{\normalfont\large\UL@phvfamily}
\fi
\ifnum\UL@ptsize=12\relax
  \newcommand*{\UL@fonttitle}{\normalfont\Large\bfseries\UL@phvfamily}
  \newcommand*{\UL@fontsubtitle}{\normalfont\large\bfseries\UL@phvfamily}
  \newcommand*{\UL@fontauthor}{\normalfont\normalsize\bfseries\UL@phvfamily}
  \newcommand*{\UL@fontprogram}{\UL@fontauthor}
  \newcommand*{\UL@fontbase}{\normalfont\normalsize\UL@phvfamily}
\fi
%    \end{macrocode}
%
% \subsubsection{Interfaces interne et externe}
%
% Définition des commandes permettant de construire les pages de titre.
% L'interface utilisateur est basée sur un ensemble de commandes
% internes. On commence par celles-ci.
%    \begin{macrocode}
\newcommand{\UL@maintitle}{}
\newcommand{\UL@subtitle}{}
\newcommand*{\UL@author}{}
\newcommand*{\UL@year}{}
\newcommand*{\UL@program}{}
\newcommand*{\UL@director}{}
\newcommand*{\UL@codirector}{}
\newcommand*{\UL@nameother}{}
\newcommand*{\UL@degreeother}{}
\newcommand*{\UL@facUL}{}
\newcommand*{\UL@facother}{}
%    \end{macrocode}
% Puis les commandes visibles pour les utilisateurs, qui redéfinissent
% les commandes internes. Voir la \autoref{sec:commandes} pour leur
% signification.
%    \begin{macrocode}
\newcommand{\titre}[1]{\renewcommand{\UL@maintitle}{#1}}
\newcommand{\soustitre}[1]{%
  \UL@hassubtitletrue
  \renewcommand{\UL@subtitle}{#1}}
\newcommand*{\auteur}[1]{\renewcommand*{\UL@author}{#1}}
\newcommand*{\annee}[1]{\renewcommand*{\UL@year}{#1}}
\newcommand*{\programme}[1]{\renewcommand*{\UL@program}{#1}}
\newcommand*{\direction}[1]{\renewcommand*{\UL@director}{#1}}
\newcommand*{\codirection}[1]{\renewcommand*{\UL@codirector}{#1}}
\newcommand*{\univcotutelle}[1]{\renewcommand*{\UL@nameother}{#1}}
\newcommand*{\gradecotutelle}[1]{\renewcommand*{\UL@degreeother}{#1}}
\newcommand*{\univbidiplomation}[1]{\renewcommand*{\UL@nameother}{#1}}
\newcommand*{\gradebidiplomation}[1]{\renewcommand*{\UL@degreeother}{#1}}
\newcommand{\faculteUL}[1]{\renewcommand*{\UL@facUL}{#1}}
\newcommand*{\faculteUdeS}[1]{\renewcommand*{\UL@facother}{#1}}
\newcommand*{\faculteUQO}[1]{\renewcommand*{\UL@facother}{#1}}
\newcommand*{\faculteUQAC}[1]{\renewcommand*{\UL@facother}{#1}}
%    \end{macrocode}
%
% \subsubsection{Titre et sous-titre}
%
% Le titre et le sous-titre peuvent s'étendre sur plus d'une ligne.
% Sans traitement spécial, un long titre ou sous-titre aurait pour
% impact de décaler vers le bas tous les autres éléments de la page
% de titre. Pour contrer ce phénomène, nous devrons mesurer la hauteur du
% titre et du sous-titre pour ensuite ajuster en conséquence la
% distance entre ce bloc et les éléments qui suivent.
%
% \begin{macro}{\UL@measuretitle}
%   On place le titre et le sous-titre centrés dans des boîtes
%   \cmd{\UL@titlebox} et \cmd{\UL@subtitlebox}. La commande
%   \cmd{\UL@measuretitle} permettra de mesurer leur hauteur lorsque le
%   titre sera créé avec \cmd{\pagestitre}, plus loin. Un espacement
%   vertical d'un demi interligne est ajouté entre le titre et le
%   sous-titre, le cas échéant.
%    \begin{macrocode}
\newsavebox{\UL@titlebox}
\newsavebox{\UL@subtitlebox}
\newlength{\UL@titleboxtotht}
\newlength{\UL@subtitleboxtotht}
\newcommand{\UL@measuretitle}{%
  \setbox\UL@titlebox=\vbox{%
    \centering\UL@fonttitle\UL@maintitle}
  \setlength{\UL@titleboxtotht}{%
    \dimexpr\ht\UL@titlebox+\dp\UL@titlebox}
  \ifUL@hassubtitle
    \setbox\UL@subtitlebox=\vbox{%
      \centering\vspace*{0.5\baselineskip}\UL@fontsubtitle\UL@subtitle}
    \setlength{\UL@subtitleboxtotht}{%
      \dimexpr\ht\UL@subtitlebox+\dp\UL@subtitlebox}
  \fi}
%    \end{macrocode}
% \end{macro}
%
% \subsubsection{Type de document}
%
% \begin{macro}{\UL@typeofdoc}
%   La commande |\UL@typeofdoc| contient le type de document qui est
%   produit: thèse, thèse en cotutelle, maîtrise
%   ou examen de doctorat.
%    \begin{macrocode}
\ifUL@isthesis
  \ifUL@iscotutelle
    \newcommand*{\UL@typeofdoc}{Th\`ese en cotutelle}
  \else
    \newcommand*{\UL@typeofdoc}{Th\`ese}
  \fi
\else
  \newcommand*{\UL@typeofdoc}{M\'emoire}
\fi
\ifUL@isexam
  \renewcommand*{\UL@typeofdoc}{Examen de doctorat}
\fi
%    \end{macrocode}
% \end{macro}
%
% \begin{macro}{\UL@docid}
%   La commande |\UL@docid| prépare ensuite la mention mise en forme
%   du type de document pour les pages de titre. La thèse ou le
%   mémoire en cotutelle ou en bidiplomation requiert un traitement
%   différent puisque le programme d'étude apparaît immédiatement sous
%   la mention.
%    \begin{macrocode}
\newcommand{\UL@docid}{%
  {\UL@fontprogram\UL@typeofdoc\par
  \ifnum\UL@typenum=2 \UL@program\par \fi}}
%    \end{macrocode}
% \end{macro}
%
% \subsubsection{Détails sur les facultés et universités d'attache}
%
% \begin{macro}{\Ul@details}
%   La commande |\Ul@details| est la plus complexe puisque la
%   disposition des informations additionnelles sur le document varie
%   beaucoup selon le type de thèse ou de mémoire. Il existe quatre
%   grandes catégories de disposition des éléments sur la page de titre:
%   standard; multifacultaire; en cotutelle ou en bidiplomation
%   (disposition identique); en extension.
%
%   Tel qu'expliqué à la \autoref{sec:utilisation:options}, certains
%   types de grade requièrent expressément que certaines informations
%   soient fournies. Si un élément d'information manque, un
%   avertissement est émis.
%    \begin{macrocode}
\newcommand{\UL@details}{%
  \ifcase\UL@typenum\relax% 0 standard
    \vspace{96pt}
    {\UL@fontprogram\UL@program}\par
    \UL@degree\par
    \vspace{112pt}
    Qu\'ebec, Canada\par
  \or%                      1 multifacultaire
    \vspace{96pt}
    {\UL@fontprogram\UL@program}\par
    \UL@degree\par
    \vspace{36pt}
    \ifx\UL@facUL\empty
      \ClassWarningNoLine{ulthese}{UL faculty names missing.}
    \else
      \UL@facUL\par
    \fi
    \vspace{48pt}
    Qu\'ebec, Canada\par
  \or%                      2 cotutelle et bidiplomation
    \vspace{72pt}
    Universit\'e Laval\par Qu\'ebec, Canada\par
    \UL@degree\par
    \vspace{\baselineskip} et\par \vspace{\baselineskip}
    \ifx\UL@nameother\empty
      \ClassWarningNoLine{ulthese}{Other university name and location missing}
    \else
      \UL@nameother\par
    \fi
    \ifx\UL@degreeother\empty
      \ClassWarningNoLine{ulthese}{Other university degree missing}
    \else
      \UL@degreeother\par
    \fi
  \or%                      3 extension
    \vspace{48pt}
    {\UL@fontprogram\UL@program\ de l'Universit\'e Laval\par
      \UL@offered\ en extension \`a l'\UL@extensionat}\par
    \vspace{36pt}
    \UL@degree\par
    \vspace{36pt}
    \ifx\UL@facother\empty
      \ClassWarningNoLine{ulthese}{Other university faculty name missing}
    \else
      \UL@facother\par
    \fi
    \UL@extensionat\par
    \UL@extensionloc\par
    \vspace{\baselineskip}
    \ifx\UL@facUL\empty
      \ClassWarningNoLine{ulthese}{UL faculty name missing}
    \else
      \UL@facUL\par
    \fi
    Universit\'e Laval\par Qu\'ebec, Canada\par
  \fi}
%    \end{macrocode}
% \end{macro}
%
% \subsubsection{Conception des pages de titre}
%
% \begin{macro}{\pagestitre}
%   Les thèses et mémoires comportent une page de titre et une page
%   frontispice. La première comporte:
%   \begin{enumerate}
%   \item le logo de l'Université Laval (sauf pour les thèses ou
%     mémoires réalisés en cotutelle, en bidiplomation ou en
%     extension);
%   \item le titre et le sous-titre, le cas échéant;
%   \item le type de document (thèse, thèse en cotutelle,
%     mémoire, etc.);
%   \item le nom complet de l'auteur;
%   \item une description du programme, du grade obtenu et des
%     facultés ou universités d'attache, le cas échéant;
%   \item la mention «Québec, Canada» si le logo de l'Université
%     Laval est présent;
%   \item la notice de copyright.
%   \end{enumerate}
%   La page frontispice reprend les éléments 2--4 et ajoute les
%   noms des directeur et codirecteurs de recherche.
%
%   On doit rétablir pour les pages de titre l'interligne simple et
%   l'espacement nul entre les paragraphes (\cmd{\parskip}). Ensuite, on
%   doit ajuster la distance entre le bloc de titre et le type de
%   document (\cmd{\UL@docidspacing}) et celle entre ce dernier et le nom
%   de l'auteur (\cmd{\UL@authorspacing}). Cela fait en sorte que les
%   éléments des pages de titre se retrouvent (presque) toujours au même
%   endroit sur la page. Une distance minimale d'un interligne est
%   conservée entre le bloc de titre et le type de document
%   (précaution nécessaire pour l'éventuel cas d'un bloc de titre
%   s'étendant sur plusieurs lignes).
%
%   Sur la page frontispice, le logo de l'Université est remplacé par
%   une boîte de réglure invisible (\emph{strut}) de la même hauteur.
%
%   Le nom de l'auteur, les directeurs et codirecteurs de recherche et
%   la notice de copyright sont insérées directement dans le code de
%   la commande \cmd{\pagestitre}.
%    \begin{macrocode}
\newlength{\UL@docidspacing}
\setlength{\UL@docidspacing}{82pt}
\newlength{\UL@authorspacing}
\setlength{\UL@authorspacing}{72pt}
\newcommand{\pagestitre}{{%
    \clearpage
    \pagestyle{empty}
    \SingleSpacing\setlength{\parskip}{0pt}
    \centering
    \UL@fontbase
    \UL@measuretitle
    \addtolength{\UL@docidspacing}{-\UL@titleboxtotht}
    \addtolength{\UL@docidspacing}{-\UL@subtitleboxtotht}
    \ifdim\UL@docidspacing<\baselineskip\relax
      \setlength{\UL@docidspacing}{\baselineskip}
      \addtolength{\UL@authorspacing}{-\baselineskip}
    \fi
    \ifnum\UL@typenum>1\relax
      \vspace*{0pt}\par
    \else
      \includegraphics[height=15mm,keepaspectratio=true]{ul_p}\par
    \fi
    \vspace{82pt}
    \copy\UL@titlebox
    \copy\UL@subtitlebox
    \vspace{\UL@docidspacing}
    \UL@docid
    \vspace{\UL@authorspacing}
    {\UL@fontauthor\UL@author}\par
    \UL@details
    \vfill
    {\textcopyright} \UL@author, \UL@year\par
    \ifUL@isexam\else
      \clearpage
      \ifnum\UL@typenum>1\relax
        \vspace*{0pt}\par
      \else
        \rule{0mm}{15mm}\par    % strut
      \fi
      \vspace{82pt}
      \box\UL@titlebox
      \box\UL@subtitlebox
      \vspace{\UL@docidspacing}
      \UL@docid
      \vspace{\UL@authorspacing}
      {\UL@fontauthor\UL@author}\par
      \vspace{72pt}
      Sous la direction de:\par
      \vspace{\baselineskip}
      \UL@director\par
      \UL@codirector
    \fi
    \clearpage}}
%    \end{macrocode}
% \end{macro}
%
% \begin{macro}{\pagetitre}
%   La page de titre était produite avec la commande \cmd{\pagetitre} dans
%   les versions précédentes de la classe. Émettre un avertissement et
%   utiliser plutôt \cmd{\pagestitre} si la commande obsolète est utilisée.
%    \begin{macrocode}
\newcommand{\pagetitre}{
  \ClassWarning{ulthese}{Command \protect\pagetitre\space is obsolete.\MessageBreak
    Using \protect\pagestitre\space instead}\pagestitre}
%    \end{macrocode}
% \end{macro}
%
% \subsection{Listes des figures et des tableaux}
%
% \begin{macro}{\listfigurename}
%   Le paquetage \pkg{babel} définit comme titre pour la liste des
%   figures «Table des figures», alors que la liste des tableaux est
%   «Liste des tableaux». Pour une plus grande symétrie, la classe
%   redéfinit le titre correspondant à \cmd{\listoffigures}. La commande
%   \cmd{\addto} est nécessaire pour éviter que \pkg{babel} redéfinisse
%   le titre à |\begin{document}|.
%    \begin{macrocode}
\ifUL@babel
  \addto\captionsfrench{\renewcommand{\listfigurename}{Liste des figures}}
\fi
%    \end{macrocode}
%   Si \pkg{babel} n'est pas chargé, ce sera à l'utilisateur de faire
%   une correction équivalente. Avec \pkg{polyglossia}, la commande à
%   insérer dans l'entête du document est la même que ci-dessus.
% \end{macro}
%
% \subsection{Dédicace et épigraphe}
%
% La dédicace et l'épigraphe sont mises en forme avec la commande
% \cmd{\epigraph} de \class{memoir}.
% \begin{macro}{\dedicace}
%   La dédicace est une épigraphe simplifiée placée seule sur une
%   page, alignée à droite à une dizaine de lignes de la marge
%   supérieure, sans auteur ou source et sans ligne de démarcation.
%    \begin{macrocode}
\newcommand{\dedicace}[1]{{%
    \clearpage
    \pagestyle{empty}
    \setlength{\beforeepigraphskip}{10\baselineskip}
    \setlength{\epigraphrule}{0pt}
    \epigraphtextposition{flushright}
    \mbox{}\epigraph{\itshape #1}{}}}
%    \end{macrocode}
% \end{macro}
% \begin{macro}{\epigraphe}
%   L'épigraphe de début de document est placée seule sur une page à
%   une dizaine de lignes de la marge supérieure. Pour le reste, on
%   s'en remet à la commande \cmd{\epigraph} de \class{memoir}.
%    \begin{macrocode}
\newcommand{\epigraphe}[2]{{%
    \clearpage
    \pagestyle{empty}
    \setlength{\beforeepigraphskip}{10\baselineskip}
    \mbox{}\epigraph{#1}{#2}}}
%    \end{macrocode}
% \end{macro}
%
% \subsection{Citations}
%
% \begin{environment}{quote}
%   La classe redéfinit l'environnement |quote| de \class{memoir}
%   afin que le texte des citations se trouve en retrait de 10~mm à
%   gauche et à droite, conformément aux règles de présentation de la
%   FESP.
%    \begin{macrocode}
\renewenvironment{quote}{%
  \list{}{\rightmargin 10mm \leftmargin 10mm}%
  \item[]}{\endlist}
%    \end{macrocode}
%   \end{environment}
%
% \begin{environment}{quotation}
%   Il en va de même de l'environnement |quotation|. Cependant, cet
%   environnement passe également à l'interligne simple et la classe
%   ajuste l'espacement vertical entre les paragraphes afin que
%   ceux-ci soient bien distincts les uns des autres tout en demeurant
%   raisonnablement compacts. Cet espacement est ici fixé à 6~points.
%    \begin{macrocode}
\renewenvironment{quotation}{%
  \list{}{%
    \SingleSpacing
    \listparindent 0em
    \itemindent    \listparindent
    \leftmargin    10mm
    \rightmargin   \leftmargin
    \parsep        6\p@ \@plus\p@}%
  \item[]}{\endlist}
%    \end{macrocode}
% \end{environment}
%
% \subsection{Numérotation des divisions du document}
%
% Par défaut, \class{memoir} numérote les divisions du document
% seulement jusqu'au niveau des sections. La classe étend la
% numérotation aux sous-sections.
%    \begin{macrocode}
\setsecnumdepth{subsection}
%</class>
%    \end{macrocode}
% ^^A Fin du code de la classe
%
% \Finale
%
% \iffalse
% ^^A Gabarits du document maître
%<*gabarit>
%<phd&standard>%% GABARIT POUR THÈSE STANDARD
%<phd&mesure>%% GABARIT POUR THÈSE SUR MESURE
%<phd&articles>%% GABARIT POUR THÈSE PAR ARTICLES
%<phd&multifac>%% GABARIT POUR THÈSE MULTIFACULTAIRE
%<phd&cotutelle>%% GABARIT POUR THÈSE EN COTUTELLE
%<phd&UdeS>%% GABARIT POUR THÈSE EN EXTENSION À L'UNIVERSITÉ DE SHERBROOKE
%<phd&UQO>%% GABARIT POUR THÈSE EN EXTENSION À L'UQO
%<m&standard>%% GABARIT POUR MÉMOIRE STANDARD
%<m&mesure>%% GABARIT POUR MÉMOIRE SUR MESURE
%<m&bidiplomation>%% GABARIT POUR MÉMOIRE EN BIDIPLOMATION
%<m&UQAC>%% GABARIT POUR MÉMOIRE EN EXTENSION À L'UQAC
%%
%% Consulter la documentation de la classe ulthese pour une
%% description détaillée de la classe, de ce gabarit et des options
%% disponibles.
%%
%% [Ne pas hésiter à supprimer les commentaires après les avoir lus.]
%%
%% Déclaration de la classe avec le type de grade
%<phd>%%   [l'un de LLD, DMus, DPsy, DThP, PhD]
%<m>%%   [l'un de MATDR, MArch, MA, LLM, MErg, MMus, MPht, MSc, MScGeogr,
%<m>%%    MServSoc, MPsEd]
%% et les langues les plus courantes. Le français sera la langue par
%<!articles>%% défaut du document.
%<articles>%% défaut du document. L'option 'bibsection' permet de créer des
%<articles>%% bibliographies par chapitre présentées sous forme de section
%<articles>%% numérotée.
%<phd&(standard|mesure)>\documentclass[PhD,english,french]{ulthese}
%<phd&articles>\documentclass[PhD,bibsection,english,french]{ulthese}
%<phd&multifac>\documentclass[PhD,multifacultaire,english,french]{ulthese}
%<phd&cotutelle>\documentclass[PhD,cotutelle,english,french]{ulthese}
%<phd&UdeS>\documentclass[PhD,extensionUdeS,english,french]{ulthese}
%<phd&UQO>\documentclass[PhD,extensionUQO,english,french]{ulthese}
%<m&(standard|mesure)>\documentclass[MSc,english,french]{ulthese}
%<m&bidiplomation>\documentclass[MA,bidiplomation,english,french]{ulthese}
%<m&UQAC>\documentclass[MSc,extensionUQAC,english,french]{ulthese}
  %% Encodage utilisé pour les caractères accentués dans les fichiers
  %% source du document. Les gabarits sont encodés en UTF-8. Inutile
  %% avec XeLaTeX, qui gère Unicode nativement.
  \ifxetex\else \usepackage[utf8]{inputenc} \fi

  %% Charger ici les autres paquetages nécessaires pour le document.
  %% Quelques exemples; décommenter au besoin.
  %\usepackage{amsmath}       % recommandé pour les mathématiques
  %\usepackage{ncccomma}      % gestion de la virgule dans les nombres

  %% Utilisation d'une autre police de caractères pour le document.
  %% - Sous LaTeX
  %\usepackage{mathpazo}      % texte et mathématiques en Palatino
  %\usepackage{mathptmx}      % texte et mathématiques en Times
  %% - Sous XeLaTeX
  %\setmainfont{TeX Gyre Pagella}      % texte en Pagella (Palatino)
  %\setmathfont{TeX Gyre Pagella Math} % mathématiques en Pagella (Palatino)
  %\setmainfont{TeX Gyre Termes}       % texte en Termes (Times)
  %\setmathfont{TeX Gyre Termes Math}  % mathématiques en Termes (Times)

  %% Gestion des hyperliens dans le document. S'assurer que hyperref
  %% est le dernier paquetage chargé.
  \usepackage{hyperref}
  \hypersetup{colorlinks,allcolors=ULlinkcolor}

  %% Options de mise en forme du mode français de babel. Consulter la
  %% documentation du paquetage babel pour les options disponibles.
  %% Désactiver (effacer ou mettre en commentaire) si l'option
  %% 'nobabel' est spécifiée au chargement de la classe.
  \frenchbsetup{%
    StandardItemizeEnv=true,       % format standard des listes
    ThinSpaceInFrenchNumbers=true, % espace fine dans les nombres
    og=«, fg=»                     % caractères « et » sont les guillemets
  }

%<!articles>  %% Style de la bibliographie.
%<!articles> \bibliographystyle{}
%<articles>  %% Suppression du numéro de section de la bibliographie. Utilisation
%<articles>  %% de \extrasfrench parce que c'est la dernière langue déclarée dans
%<articles>  %% \documentclass, ci-dessus.
%<articles>  %\addto\extrasfrench{%
%<articles>  %  \renewcommand{\bibsection}{\section*{\bibname}\prebibhook}}

  %% Déclarations des pages de titre. Remplacer les éléments entre < >.
  %% Supprimer les caractères < >. Couper un long titre ou un long
  %% sous-titre manuellement avec \\.
  \titre{<Titre principal>}
  % \titre{Ceci est un exemple de long titre \\
  %   avec saut de ligne manuel}
  % \soustitre{Sous-titre le cas échéant}
  % \soustitre{Ceci est un exemple de long sous-titre \\
  %   avec saut de ligne manuel}
  \auteur{<Prénom Nom>}
  \annee{<20xx>}
%<phd&(standard|multifac|cotutelle)>  \programme{Doctorat en <discipline> <-- majeure, s'il y a lieu>}
%<phd&(mesure)>  \programme{Doctorat sur mesure en <discipline> <-- majeure, s'il y a lieu>}
%<phd&(UdeS|UQO)>  \programme{Doctorat en <discipline>}
%<m&standard>  \programme{Maîtrise en <discipline> <-- majeure, s'il y a lieu>}
%<m&mesure>  \programme{Maîtrise sur mesure en <discipline> <-- majeure, s'il y a lieu>}
%<m&bidiplomation>  \programme{Maîtrise en <discipline>}
%<m&UQAC>  \programme{Maîtrise en <discipline>}
%<!(cotutelle|bidiplomation)>  \direction{<Prénom Nom>, <directeur ou directrice> de recherche}
%<cotutelle>  \direction{<Prénom Nom>, <directeur ou directrice> de recherche \\
%<cotutelle>             <Prénom Nom>, <directeur ou directrice> de cotutelle}
%<bidiplomation>  \direction{<Prénom Nom>, <directeur ou directrice> de recherche \\
%<bidiplomation>             <Prénom Nom>, <directeur ou directrice> de recherche}
  % \codirection{<Prénom Nom>, <codirecteur ou codirectrice> de recherche}
  % \codirection{<Prénom Nom>, <codirecteur ou codirectrice> de recherche \\
  %              <Prénom Nom>, <codirecteur ou codirectrice> de recherche}
%<cotutelle>  \univcotutelle{<Université de cotutelle> \\ <Ville>, <Pays>}
%<bidiplomation>  \univcotutelle{<Université de bidiplomation> \\ <Ville>, <Pays>}
%<cotutelle|bidiplomation>  \gradecotutelle{<Nom du grade> (<sigle du grade>)}
%<multifac>  \faculteUL{<Faculté 1> \\ <Faculté 2>}
%<UdeS|UQO|UQAC>  \faculteUL{<Nom de la faculté UL>}
%<UdeS>  \faculteUdeS{<Nom de la faculté UdeS>}
%<UQO>  \faculteUQO{<Nom de la faculté UQO>}
%<UQAC>  \faculteUQAC{<Nom de la faculté UQAC>}

\begin{document}

\frontmatter                    % pages liminaires

\pagestitre                     % production des pages de titre

\chapter*{Résumé}                      % ne pas numéroter
\phantomsection\addcontentsline{toc}{chapter}{Résumé} % inclure dans TdM

\begin{otherlanguage*}{french}
  Texte du résumé en français.
\end{otherlanguage*}
                % résumé français
\chapter*{Abstract}                      % ne pas numéroter
\phantomsection\addcontentsline{toc}{chapter}{Abstract} % inclure dans TdM

\begin{otherlanguage*}{english}
  Text of English abstract.
\end{otherlanguage*}
              % résumé anglais
\cleardoublepage

\tableofcontents                % production de la TdM
\cleardoublepage

\listoftables                   % production de la liste des tableaux
\cleardoublepage

\listoffigures                  % production de la liste des figures
\cleardoublepage

\dedicace{Dédicace si désiré}
\cleardoublepage

\epigraphe{Texte de l'épigraphe}{Source ou auteur}
\cleardoublepage

\chapter*{Remerciements}         % ne pas numéroter
\phantomsection\addcontentsline{toc}{chapter}{Remerciements} % inclure dans TdM

Texte des remerciements en prose.
         % remerciements
\chapter*{Avant-propos}         % ne pas numéroter
\phantomsection\addcontentsline{toc}{chapter}{Avant-propos} % inclure dans TdM

L'avant-propos est surtout nécessaire pour une thèse par article.
           % avant-propos

\mainmatter                     % corps du document

\chapter*{Introduction}         % ne pas numéroter
\phantomsection\addcontentsline{toc}{chapter}{Introduction} % inclure dans TdM

Une thèse ou un mémoire devrait normalement débuter par une
introduction. Celle-ci est traitée comme un chapitre normal, sauf
qu'elle n'est pas numérotée.
          % introduction
%<!articles>\chapter{Architecture du robot 3DDL}     % numéroté

Texte du chapitre.
             % chapitre 1
%<articles>\chapter{Titre du chapitre}     % numéroté

Texte du chapitre.

\bibliographystyle{}              % style de la bibliographie
\bibliography{}                   % production de la bibliographie
    % chapitre 1
%<!articles>\chapter{Développement de l'environnement de développement, robot 3DDL}     % numéroté

Texte du chapitre.
             % chapitre 2, etc.
%<articles>\chapter{Titre du chapitre}     % numéroté

Texte du chapitre.

\bibliographystyle{}              % style de la bibliographie
\bibliography{}                   % production de la bibliographie
    % chapitre 2, etc.
\chapter*{Conclusion}         % ne pas numéroter
\phantomsection\addcontentsline{toc}{chapter}{Conclusion} % dans TdM

Une thèse ou un mémoire devrait normalement se terminer par une
conclusion, placée avant les annexes, le cas échéant. Celle-ci est
traitée comme un chapitre normal, sauf qu'elle n'est pas numérotée.
            % conclusion

\appendix                       % annexes le cas échéant

\chapter{Titre de l'annexe}     % numérotée

Texte de l'annexe.
                % annexe A

%<!articles>\bibliography{}                 % production de la bibliographie
%<!articles>
\end{document}
%</gabarit>
%
% ^^A Gabarits des parties du document
%<*resume>
\chapter*{Résumé}                      % ne pas numéroter
\phantomsection\addcontentsline{toc}{chapter}{Résumé} % inclure dans TdM

\begin{otherlanguage*}{french}
  Texte du résumé en français.
\end{otherlanguage*}
%</resume>
%
%<*abstract>
\chapter*{Abstract}                      % ne pas numéroter
\phantomsection\addcontentsline{toc}{chapter}{Abstract} % inclure dans TdM

\begin{otherlanguage*}{english}
  Text of English abstract.
\end{otherlanguage*}
%</abstract>
%
%<*remerciements>
\chapter*{Remerciements}         % ne pas numéroter
\phantomsection\addcontentsline{toc}{chapter}{Remerciements} % inclure dans TdM

Texte des remerciements en prose.
%</remerciements>
%
%<*avantpropos>
\chapter*{Avant-propos}         % ne pas numéroter
\phantomsection\addcontentsline{toc}{chapter}{Avant-propos} % inclure dans TdM

L'avant-propos est surtout nécessaire pour une thèse par article.
%</avantpropos>
%
%<*introduction>
\chapter*{Introduction}         % ne pas numéroter
\phantomsection\addcontentsline{toc}{chapter}{Introduction} % inclure dans TdM

Une thèse ou un mémoire devrait normalement débuter par une
introduction. Celle-ci est traitée comme un chapitre normal, sauf
qu'elle n'est pas numérotée.
%</introduction>
%
%<*chapitre>
\chapter{Titre du chapitre}     % numéroté

Texte du chapitre.
%<articles>
%<articles>\bibliographystyle{}              % style de la bibliographie
%<articles>\bibliography{}                   % production de la bibliographie
%</chapitre>
%
%<*conclusion>
\chapter*{Conclusion}         % ne pas numéroter
\phantomsection\addcontentsline{toc}{chapter}{Conclusion} % dans TdM

Une thèse ou un mémoire devrait normalement se terminer par une
conclusion, placée avant les annexes, le cas échéant. Celle-ci est
traitée comme un chapitre normal, sauf qu'elle n'est pas numérotée.
%</conclusion>
%
%<*annexe>
\chapter{Titre de l'annexe}     % numérotée

Texte de l'annexe.
%</annexe>
% Local Variables:
% mode: doctex
% coding: utf-8
% TeX-master: t
% TeX-engine: xetex
% End:
% \fi
